% =====================================================================
%  R. Nielsen, OM PHANTASIENS MAGT.  Sex Forelæsninger.
%  Kjøbenhavn: Forlagt af den Gyldendalske Boghandel (F. Hegel); trykt hos J. Jørgensen & Co., 1876.  91 pp.  Antiqua.
%
%  TRANSCRIPTION -- GENERATED; DO NOT EDIT BY HAND.  Rebuild from the repo root:
%    python3 ebook-method/phantasiens-magt/draft.py --emit --force \
%      && python3 ebook-method/collate/check.py phantasiens-magt --italics --apply \
%      && python3 ebook-method/phantasiens-magt/draft.py --post
%
%  STATUS: COMPLETE (2026-10-12): title page, Indhold and the whole text, pp. 1-91.  The book has no preface and no notes.
%
%  EDITION.  The text of record is the Royal Library's scan (KB 130025470861; ~/bibliotek/Nielsen, Rasmus/
%  phantasiens-magt.pdf, SCANS.tsv), 300 ppi colour, with KB's OCR as text layer.  The only printing; no e-book and no
%  earlier transcription exists (checked 2026-10-12: the repo, the web, KB).  This copy is marked in pencil throughout
%  (underlinings, strokes and brackets in the margins), which the layer read as letters and which split words.
%
%  ACCURACY (2026-10-12).  The words are those of KB's layer (draft.py), collated word by word with an independent
%  tesseract reading of the images (collate/check.py, witness 'ocr'): 569 disagreements in 23,083 words.  About 80 were
%  settled by check.py's own rules and 104 by rules.py, each rule grounded in this book's crops: the print sets „...“ only
%  (tesseract's » and % are its reading of „ and “: zoom q-01); tesseract's aa read as an/ag (maae>mane, faae>fane);
%  tesseract's glyph faults where the layer's word is attested and tesseract's not; the layer's glyph faults of equal
%  length where tesseract's word is attested and the layer's not; tesseract's marks where the layer has nothing and the
%  print has pencil (this copy's); accents tesseract adds; the printer's space before ;:?!.  The candidate rules were
%  measured first on a blind pilot of 58 crops (8/8, 12/12, 8/8, and 14/15 -- the miss, „fiint“, which both OCRs read
%  wrong, led to the equal-length restriction).  The rest, 385 checks in all with the pilot, were read at the image BLIND
%  (sheets.py: crops with no readings shown, two planted controls a sheet, 78 of 78 controls read right; six Sonnet
%  readers, .parts/reads/r-b1-A.tsv, r-b2-B..F.tsv); about 110 of those were decided in the main session from the
%  readers' lines, and some 20 spots at 300-600 dpi (zoom.py, sheets q, s, t, u, v, w, x600, y600, z).
%  Then four sweeps of the collated text, each spot read at the image: check.py --splits and a corpus-vocabulary sweep
%  for words the layer split where pencil underlining crosses them (18 mended: k unn e, e gen Erkj endelseskreds ...);
%  check.py --moved for widely spaced lines the layer put out of order (pp. 21, 60, 66, 83, 85, 87, 89, 91); the readers'
%  whole LINEs matched against the text (found „kan Tro sættes“, p. 44, where a W decision had lost „Tro“, and
%  „Feltraabet“, p. 67, which both OCRs read „Feltraahet“); and every word in no other transcription of the corpus (383
%  words; found the misprints below, a stain and two specks both OCRs read as letters).  No italics (slant scores
%  99 % <= 1.04, no runs) and no letterspacing in the body (the twelve highest letter-gap scores were all unspaced); the
%  Indhold letterspaces »Forelæsning«.  SAMPLED PAGE CHECK (pagecheck.py, seed 202610121): pp. 19, 28, 84 and the title
%  page read whole and blind, 778 words: 0 word, 0 punctuation and 0 markup errors (95 % upper bound 3.9 per 1000 words).
%  What this cannot see: a misprint both OCRs normalise into a real word -- the corpus sweep finds only non-words.
%  TRANSLATION as a last reading: the four translators listed 40 odd spots (ebook-method/phantasiens-magt/tr/odd-*.tsv),
%  all read at the image (zoom o-01, o-02, q62): 16 mended by POST patch -- a line-end word the layer lost (p. 41
%  „Phantasie.“), four false paragraph breaks from pencil in the margin (pp. 29, 60), four raised specks read as stops or
%  hyphens by both OCRs (pp. 11, 31, 48, 62), „aftale“ for „at tale“ (p. 67) -- and four more misprints found; the rest
%  are the print's (Øjne, nok som, Rowel, sindig, Vexelskin, quotation marks the print leaves open).
%  PRINTED AS IS (each seen at the image; a % comment at the site): „Tankespørgsmal“ and „Nodern“ (p. 14), „kar“ for
%  har (p. 55), „et“ for er (p. 61), „sil“ for til (p. 63), „Chisti“ (p. 68), „maaksee“ (p. 69), „ekstastiske“ (p. 88);
%  quotations not closed in the print (pp. 56, 59).  Specks read as accents by both OCRs: „Herregaarden“ (p. 2), „havde“ (p. 61).
%  The ø rule (Hans 2026-10-08) found nothing here (oefix.py's „stode“ -> „støde“ is wrong: stode, „stood“, is a word).
%  Not transcribed: signature marks („1“, „1*“ ...), library stamps, the copy's pencil marks; the rule under each lecture
%  heading is drawn by \lecture.  Quotation marks are as printed („ “).
%
%  PAGE MAP.  printed p = PDF page - 9 (p. 1 = PDF 10, p. 91 = PDF 100); title PDF 6, Indhold PDF 8.
%  Each lecture begins on a new page: I p. 1, II p. 14, III p. 28, IV p. 43, V p. 57, VI p. 75.
%  The Indhold sets »Første Forelæsning« etc. letterspaced; reproduced here in \emph (house convention).
%
%  TOOLS (ebook-method/phantasiens-magt/, copied from betingelserne-villie): draft.py, ocr.py, rules.py, sheets.py +
%  reads.py + READER-BRIEF.md, zoom.py, prescreen.py, pagecheck.py + PAGECHECK-BRIEF.md, patches.py (PATCHES, POST),
%  decisions.tsv.
%
% =====================================================================
\documentclass[12pt,a4paper,openany]{book}
\usepackage[utf8]{inputenc}
\usepackage[T1]{fontenc}
\usepackage{libertinus}
\usepackage{libertinust1math}
\usepackage[danish]{babel}
\usepackage[protrusion=true,expansion=true]{microtype}
\usepackage{setspace}
\usepackage[margin=1.3in]{geometry}
\usepackage{fancyhdr}
\usepackage{hyperref}
\hypersetup{pdftitle={Om Phantasiens Magt}, pdfauthor={Rasmus Nielsen}, hidelinks}
\pagestyle{fancy}
\fancyhf{}
\fancyhead[LE,RO]{\thepage}
\fancyhead[C]{\textit{R. Nielsen: Om Phantasiens Magt (1876)}}
\setstretch{1.3}
\usepackage{marginnote}
\renewcommand{\thefootnote}{\arabic{footnote})}
\newcommand{\tocline}[3]{\noindent\makebox[2.2em][r]{#1}\hspace{0.6em}#2\dotfill\ #3\par}
\newcommand{\opage}[1]{\marginnote{\footnotesize\textit{[#1]}}}
% a lecture opens a new page: heading, a short rule; the print's large initial is not reproduced
\newcommand{\lecture}[1]{\clearpage\begin{center}{\large #1}\\[0.4em]\rule{3em}{0.4pt}\end{center}\medskip\noindent}

\begin{document}

% --- front matter: PDF 6, title (read at the image, 2026-10-12); the verso carries nothing of the text.
\begin{titlepage}
\begin{center}
\vspace*{3em}
{\large OM}\\[1em]
{\Huge PHANTASIENS MAGT.}\\[2.5em]
SEX FORELÆSNINGER\\[1em]
{\small AF}\\[1em]
{\large R. NIELSEN.}\\[2em]
[publisher's device]\\[4em]
KJØBENHAVN.\\[0.5em]
FORLAGT AF DEN GYLDENDALSKE BOGHANDEL (F. HEGEL).\\[0.5em]
{\small TRYKT HOS J. JØRGENSEN \& CO.}\\[0.5em]
{\small 1876.}
\end{center}
\end{titlepage}

% --- PDF 8: Indhold (read at the image, 2026-10-12)
\begin{center}{\large Indhold.}\\[0.3em]\rule{3em}{0.4pt}\end{center}
\hfill{\small Side}\par
\noindent\emph{Første Forelæsning.}\enspace Phantasie og Virkelighed\dotfill\ 1.\par
\noindent\emph{Anden Forelæsning.}\enspace Phantasiens Forhold til Spørgsmaalet\par
\noindent\hspace*{2em}om Sjælens Udødelighed\dotfill\ 14.\par
\noindent\emph{Tredie Forelæsning.}\enspace Religiøs Phantasie\dotfill\ 28.\par
\noindent\emph{Fjerde Forelæsning.}\enspace Gammeltestamentlig Phantasie\dotfill\ 43.\par
\noindent\emph{Femte Forelæsning.}\enspace Nytestamentlig Phantasie\dotfill\ 57.\par
\noindent\emph{Sjette Forelæsning.}\enspace Phantasiens uendelige Rige\dotfill\ 75.\par
\clearpage


% ===== TEXT (pp. 1--91) =====
% --- p. 1 ---
\setcounter{footnote}{0}%
\lecture{Første Forelæsning.}\opage{1}

Vi kjende det daglige Liv, dets Glæder og Sorger, dets
Hygge og Fred. Alt er os her saa hjemligt, saa tilvant,
saa bekjendt. Men bagved det Bekjendte er der et
Ubekjendt, og dette Ubekjendte melder sig stundom
uvilkaarlig med Spørgsmaal, som ængste. Vi sidde i
Familiens Kreds i Solskin af huuslig Lykke, det Ubekjendte
hvisker: hvorlænge vil det vare? Vi mindes
vore Hedengangne, og Religionen trøster os med et
Gjensyn; det Ubekjendte hvisker: er der et Gjensyn,
er Du vis derpaa? Saaledes i mangfoldige Tilfælde.
Hvad er nu dette Ubekjendte? Er det af det Onde
eller af det Gode? Hvorfra stammer det, fra Lysets
Ophav eller fra Mørkets Egne? Det kommer som
et Iilbud, et Varsel, en Mindelse om det Usynlige,
en manende Røst fra Aandernes Verden. Der er
Mennesker, som have en naturlig Sky for at gjøre dette
Ubekjendtes nærmere Bekjendtskab. Men Hverdagslivet
selv er underlagt Tilskikkelser, som kunne bringe det
Ubekjendte til at træde frem, pludselig og i al sin Vælde.
Lad mig anføre et lille Exempel.

Et ungt Menneske, med oprindelige Anlæg og bestemt
til Studeringer, havde allerede gjennemgaaet det
% --- p. 2 ---
\setcounter{footnote}{0}%
\opage{2}første akademiske Cursus, da Faderen, som begyndte
at føle sig gammel og træt, ønskede at hans Søn skulde
blive Landmand og hjælpe ham med at bestyre Gaarden.
Sønnen gav efter for Faderens Ønske, saa meget villigere,
som der i Egnen var en Herregaard, og paa Herregaarden % scan: a speck over the second a
en ung Frøken, Godseierens Datter. Den unge Mand
havde seet paa hende, og hun havde igjen seet paa ham ---
hans Beslutning var tagen. Hen i Høst --- han havde
allerede den Gang været et Par Aar i sin nye Virksomhed
--- viste sig en Komet paa Himlen. Det gav
Folk Noget at tænke paa og satte Landboernes Sind i
Røre. Gamle Koner forsikrede, at Krig, Hungersnød,
Jordskjælv og Pestilense forestode, enten alle fire Dele
eller i det Mindste een af Delene, det kunde aldrig
feile. De Yngre meente, at Sagen vel ikke var saa
farlig; Kometen var jo ordenlig anmeldt i Almanakken.
De læste og saae, at den Uge efter Uge punktlig og
stille rettede sig efter Almanakkens Forskrifter: den
kunde umulig have noget Ondt i Sinde. Men alt som
den ugeviis kom nærmere, tog Forundringen til, thi
Kometen syntes at forandre sig og antage et truende
Udseende. Det røde Øie blev bestandig hvassere, Halen
voxede, bredte og grenede sig som et stort Ildriis.
Stjernerne paa den klare Septemberhimmel fik et
mattere Skin; det var som vilde Himlen halvt lukke
Øjnene til, for ikke at see det uhyggelige Tegn. De
gamle Koner grebes af en prophetisk Aand; de prædikede
Bod og Omvendelse, og Mange bleve andægtige.

Dog det naturlige Liv gjør sine Rettigheder gjældende.
Høsten havde været velsignet; for Hungersnød
var der da ingen Fare. Rønnebærrene bleve røde, Løvet
i Skoven blev bruunt og violet og guult, og Bladene
% --- p. 3 ---
\setcounter{footnote}{0}%
\opage{3}faldt; Vinden foer rigtignok hæsblæsende hen over
Stubber og pløiede Marker, men til Jordskjælv var der
Intet at mærke; Sundhedstilstanden var god, og Aviserne
turde nok staae inde for Freden: saa skulde der da
være Høstgilde.

Høstgildet paa Herregaarden var Omegnens aarlige
og meest glimrende Fest, det var i Sandhed en folkelig
Fest; thi den liberale Godseier indbød Høie og Lave.
Piger og Fruer, Adelsmænd og Bønder blandede sig i
Dandsen mellem hverandre; Strømmen bølgede i de
store Sale og gjennem den med Fakler og Lamper oplyste
Have. Mange havde glædet sig til denne Fest,
men Ingen mere end vor unge Ven. Han ventede nu
saa vist, at denne Aften skulde krone hans Lykke.
Som han red paa Veien fattede han et Forsæt, der
skulde staae fast; han vilde denne Aften erklære sin
hemmelige Forlovelse for Godseieren; han havde omhyggelig
valgt sine Ord og saae allerede i Aanden,
hvorledes den fornemme, men humane Herre, der altid
havde paaskjønnet hans Driftighed og viist ham saa
megen Forekommenhed, lod de første Overraskelsens
Rynker glatte sig og gav efter for hans og den elskede
Datters forenede Bønner. Det forekom den unge Mand,
alt som han nærmede sig og hørte Fanfaren og saae
de vaiende Flag, at en saadan Høstfest nødvendig maatte
være bestemt til at blive en Forlovelsesfest. Forestillingen
herom blev ham saa levende, at han ordenlig begyndte
at overveie, hvorledes han mest passende skulde svare,
naar de Forlovedes Skaal blev udbragt. Men just som
han red ind ad Porten, stod det Ubekjendte paa Luur
og spurgte ham hviskende: er Du vis i din Sag?

% --- p. 4 ---
\setcounter{footnote}{0}%
\opage{4}Hvorledes Varslet slog ham, de nærmere Omstændigheder,
hvor mange piinlige Øieblikke han maatte
udholde inden han naaede Visheden, ved slige Enkeltheder
ville vi ikke opholde os. Langt om længe blev
han i een af Havens dunkle Kroge indhentet af den
unge Frøken; der var ikke Tid til Forklaring, Mængden
omsværmede dem, det raslede i Grenene, der var Folk
i Nærheden: kun dette fik hun med graadkvalt Stemme
og et Udtryk af den dybeste Smerte udtalt: Alt er forbi!
med disse Ord svandt hun som et Syn. --- „Alt forbi!“
Om en Afgrund havde aabnet sig, om Kometen havde
dasket Jorden med sin Hale, kunde den unge Mand
ikke være blevet mere bestyrtet.

„Alt forbi.“ Som en Fredløs, der ængstelig undviger
Menneskers Selskab, ilede han ad en lønlig Sti gjennem
Krattet over en Steentrappe ud paa den aabne Mark.
Stemmerne og den glade Latter tabte sig, Musikkens
Toner løde stedse svagere og døde sagte hen i den maaneklare
Nat. Alt var stille, kun enkelte fjerne Hundeglam
lode sig høre fra de omliggende Landsbyer. Han vadede
i Muld, bestandig fremad, som havde han et Maal, indtil
han standsede ved en stor Mergelgrav. Han maalte
den med Øinene; Vandet, det vidste han, var tre Gange
saa dybt som den steile Rand var høi, og denne overgik
dog en Mands Høide. Det Ubekjendte hviskede:
„Du maa bort, Du maa reise.“ Han saae i det speilede
Dyb; Kometen stod dernede, som om den holdt stille
og ventede paa ham. Det er en Paradiisfugl med gyldne
Fjer, med den vil jeg reise. Opad kunde jeg ikke naae
den, den sad mig for høit, men nedad, et Spring, et
dristigt Spring i Dybet, og jeg rider paa dens Ryg.
Men hvad er det Sorte derhenne? Kroppen af en død
% --- p. 5 ---
\setcounter{footnote}{0}%
\opage{5}Hund! En Reisekammerat, der i Hastværk har ladt
sin Overkjole i Stikken: saa kommer jeg da ikke til
at gjøre Touren alene. „Du maa først tage Afsked“,
hviskede det Ubekjendte; „Reisen er lang, hvo veed, naar
Du kommer tilbage“.

Og han vilde tage Afsked, ikke med nogen levende
Modersjæl, men med den dunkle Lund paa Høien histhenne.
Der oppe var en Græsbænk; der sov i Græsset
mellem visne Straa en Udsæd uden Høst, en Førstegrøde
af et Hjertes Vaar, forladte Børn af et skuffet
Haab, Erindringer om en kortvarig, men ubeskrivelig
Lyksalighed. Altsaa, tænkte han, først i Lunden et
Farvel, og saa tilbage til Mergelgraven og Kometen og
den døde Hund! Han ilede til Høien, der var Feber i
hans Blod, hans Hjerne brændte; han tørnede mod et
ungt men udgaaet Træ; en Ugle skreg, en opskræmmet
Hare foer ham tæt forbi. Alt løb rundt, han havde
en Fornemmelse, som om der regnede Ild fra Kometen
ned paa hans Hoved. Med vaklende Knæ arbeidede
han sig op ad Fodstien og fik Øie paa Bænken. Der
sad en kvindelig Skikkelse. Er det muligt? hun sidder
deroppe, hun venter paa mig? Nei, det er umuligt.
Jo, det er hende, og dog --- ikke hende. Den Siddende
reiste sig, belyst af Maanens Straaler, og Skikkelsen
blev aandeagtig majestætisk. „Du drager bort herfra,
vil tage Afsked; men Du skal ikke reise med den døde
Hund. En Ven skal følge Dig, en Ven, som kjender
Veien, en ædel, tro Ledsager, paalidelig som hiin, der
fulgte Beatrices Elsker igjennem Helvede, igjennem
Skjærsilden og igjennem Paradiset ind i Himlen“. Saa
svandt det underlige Syn; den unge Mand var ramt
som af et Lyn og styrtede til Jorden. Hvad der siden
% --- p. 6 ---
\setcounter{footnote}{0}%
\opage{6}blev af ham, da han reiste sig igjen, bekymrer os ikke.
Vi have kun berørt hans Historie for det Ubekjendtes
d. e. for Phantasiens Skyld; thi det Væsen, der vaagnede
til Liv i ham, da han opvaagnede, det var i Sandhed
Phantasien.

For at komme til Indsigt i Phantasiens Væsen,
maae vi helst begynde med Virkeligheden, dens Betingelser
og Love. Virkeligheden ligefrem som den staaer
og gaaer er ikke Phantasiens Gjenstand; at beskjæftige
sig med det Virkelige overlader den til Forstanden og
de fem Sandser. Men Virkeligheden skjuler Noget, som
i høi Grad interesserer Phantasien, Noget, som den
overalt udspeider og opdager; og hvad er det saa?
Spørge vi Phantasien, saa svarer den os, at det er
hemmelighedsfulde Væsener, usynlige Skjønheder, Ansigter,
som titte frem bag Virkelighedens Gardiner. Det
er da med andre Ord ikke Virkeligheden selv, men
dens Muligheder.

Men Mulighederne, de Muligheder, som sætte Phantasien
i Bevægelse, ere af høist forskjellig Art. Vi kunne
see det af vor Fortælling. De gamle Koners Phantasie
søgte Mulighederne i Kometen, men løb vild i Lovene
og miskjendte Virkeligheden. Det var astrologiske
Muligheder, og Astrologernes gyldne Tid er forbi. Ingen
troer paa Stjernernes Indflydelse, naar Spørgsmaalet er
om det jordiske Liv, dets Skjæbner og Tilstande. Tycho
Brahes Dage ere ude af Cours som forældede Mønter.
Og hvorfor? Fordi Stjernernes Stillinger, Bevægelser og
indbyrdes Vexelvirkninger rette sig efter Love, som kunne
udfindes ved Beregning. Anderledes med den Slags
Muligheder, som bedaarede den unge Mands Phantasie,
de lade sig ikke beregne, som man beregner en Komet%
% --- p. 7 ---
\setcounter{footnote}{0}%
\opage{7}bane; det er psychologiske Muligheder, Frihedsmuligheder.


Stjernerne paa den udvortes Himmel følge saa ufravigelig
bestemte Baner, at Iagttageren forud nøiagtig
kan angive det Sted, hvortil de til ethvert fastsat Tidspunkt
ville komme. Men Stjernerne paa den indvortes
Himmel, der hvor hvert enkelt Menneske helst maa lede
efter sin Lykkestjerne, bevæge sig i uregelmæssige Baner.
Vel ere ogsaa disse underkastede Love og naturlige Betingelser,
men Omstændighederne vexle, og Perturbationerne
ere saa mangfoldige, at ingen Beregning slaaer til:
Forstanden arbeider sig træt, og Phantasien leer ad Klogskabens
Forlegenhed. Vor unge Mand gik dristig frem, han
troede paa sin Lykkestjerne og blev skuffet; han kunde
lige saa godt være blevet skuffet paa aldeles modsat
Maade, ved Tilbageholdenhed og Mistro.

Blandt de Diapsalmata, der indlede S. Kierkegaards
„Enten—Eller“ forekommer følgende --- i hvad han kalder
et extatisk Foredrag, et Foredrag, som vor unge Ven
maaskee ikke har havt Tid til at høre: „Tro en Pige,
Du vil fortryde det; tro hende ikke, Du vil ogsaa fortryde det; tro en Pige eller tro hende ikke, Du vil fortryde
begge Dele; enten Du troer en Pige eller Du
troer hende ikke, Du vil fortryde begge Dele.“ Det er
et Schema, der kan udfyldes paa de forskjelligste Maader.
Fra Ethiken f. Ex.: Vær ærlig, Du vil fortryde det, vær
uærlig, Du vil ogsaa fortryde det; enten Du er ærlig
eller Du er uærlig, Du vil fortryde begge Dele. Fra
det Intellectuelle: Brug din Forstand, Du vil fortryde
det, brug den ikke, Du vil ogsaa fortryde det; brug
din Forstand eller brug den ikke, Du vil fortryde begge
Dele; enten Du bruger din Forstand eller Du ikke
% --- p. 8 ---
\setcounter{footnote}{0}%
\opage{8}bruger den, Du vil fortryde begge Dele. Saaledes kan
man jo blive ved i det Uendelige. Hvad skal nu alt
dette sige? Forfatteren giver en kort, men fyndig Forklaring:
„Hæng Dig, Du vil fortryde det; hæng Dig
ikke, Du vil ogsaa fortryde det; hæng Dig eller hæng
Dig ikke, Du vil fortryde begge Dele; enten Du hænger
Dig eller Du ikke hænger Dig, Du vil fortryde begge
Dele;“ og saa tilføier han: „Dette, mine Herrer, er Indbegrebet
af al Leve-Viisdom.“

Den Fortvivlelsens Ironie, som saaledes driver sit
dialektiske Spil, kan belyses fra mange forskjellige Sider;
en Hovedside, som vi i denne Sammenhæng ville fremhæve,
angaaer Forholdet mellem Virkelighed og Phantasie.
At lade Phantasien raade for Virkeligheden gaaer
ikke an; den vilde slaae sig løs, vende op og ned
paa Alt og rumstere forfærdeligt med den givne Tingenes
Orden. Men at udelukke Phantasien gaaer heller
ikke an. Naar det, vi i endelig Forstand kalde det
Virkelige, skal være vort Et og Alt, saa mistydes Frihedens
Muligheder, Inderlighedens Skatte, som ligge
gjemte i Sindet, forblive ubrugte, og Livskilderne udtørres.
Lad os see paa det tørre Forstandsmenneske. Hvorledes
bærer en Saadan sig ad? For ikke at bedaares af Illusioner
begynder han med at udstykke Mulighederne i
lutter forestillede Tilfælde og beregne de sandsynlige
Følger. Det kan nu være meget godt, men hvor langt
naaer han? Han opstiller to Tilfælde, tre Tilfælde,
maaskee ti, tyve forskjellige Tilfælde. Nu skal han til
at vælge. Men naar han saa vil til at vælge det ene,
falder det ham ind, at han lige saa godt kunde vælge
det modsatte; vil han holde paa det ene, saa indstiller
sig strax et andet. Han staaer omsværmet af Tilfælde,
% --- p. 9 ---
\setcounter{footnote}{0}%
\opage{9}som En, der er omsværmet af Fluer; der er noget Drillende
deri; det er Phantasiens Muligheder, der lege
Tagfat imellem Tilfældene, de have ham til Bedste, de
gjække ham lige som Alferne gjække Mester Grimmemann
ved at tilliste sig snart hans Bestalling og snart
hans Briller.

At udelukke Phantasien for at gribe Virkeligheden,
gaaer ikke an; thi man mister det Øie, hvormed man
i dybere Forstand skulde opdage Mulighederne.

Men Sagen faaer et endnu langt alvorligere Udseende,
naar Spørgsmaalet er om Menneskets Bestemmelse,
dets Stilling og Betydning i den store synlige
Verden. Bag det Synlige ligger det Usynlige; bag det
Sandselige det Usandselige; til alt det Ydre svarer der
et Indre, og naar vi see nøiere til, er det altid det
Indre der bestemmer sit Ydre. Det er jo netop det
Indre, vi kalde Tingenes Væsen, deres Anlæg, deres
Natur. Naturlæren, den forskende Erfaringslære, undersøger
Tingene i Naturen, deres Egenskaber, Kræfter,
Love; den er forvisset om, at enhver Ting maa virke i
Overeensstemmelse med sin Natur; men hvad Naturen
er, det veed den ikke. Den lægger ikke Dølgsmaal
derpaa, den skjelner skarpt mellem det, den veed, og
det, den ikke veed.

Naturlæren gaaer med denne sin Skjelnen endnu
videre; naar der handles om det Enkelte eller om hele
Grupper eller Systemer af ligeartede Kjendsgjerninger,
saa forstaaer den, at Indsigtens og Kundskabens Omraade
bestandig kan udvides. Der er Meget vi endnu
for Øieblikket ikke vide; men vi kunne med Tiden faae
det at vide. Naar der derimod spørges om det første
Oprindelige, om Tingenes Ophav, om Naturen selv, saa
% --- p. 10 ---
\setcounter{footnote}{0}%
\opage{10}standser den ved en uoverskridelig Grændse; saa tilstaaer
den besindige og ædruelige Erfaringslære, at den fra
sit Stade umulig kan ved fortsatte Undersøgelser naae
til en Viden om, hvad Tingene ere i deres inderste Grund.
„Vi kjende Tingenes Egenskaber, men ikke Tingene;
Phænomenerne, men ikke deres Væsen; vi vide ikke,
hvad Kraft er, ikke, hvad Magnetisme er, hvad Elektricitet
er, hvad Livet er og hvorfra det oprinder, Men dette
vide vi, at det Hele udvikler sig betinget i nødvendig
Sammenhæng, og at Alt maa gaae efter uforanderlige
Love“.

Dette vilde nu altsammen være saare skjønt, oplysende
og veiledende, dersom vi kunde blive staaende
herved. Men der er en Hage ved Tingen, og en Braad
i Bevidstheden. Vi tale om en indre, usynlig Verden,
der aabenbarer sig i den ydre og sandselige. Fra det,
vi see, slutte vi til det, vi ikke see. Hvorledes forholder
det sig nu med denne indre usynlige, usandselige Verden?
Gaaer alt det Indre nødvendig op i sit Ydre eller er
her et Overskud? Er det Ydre, hvori det Indre træder
frem, indskrænket til denne Verden, denne Natur, eller
er der en høiere Verden, en dybere Natur, af hvis Love,
Liv og Kræfter denne vor sandselige, synlige Verden kun
giver en ufuldkommen Afspeiling: kan der, med andre
Ord, gives en overnaturlig Verden? Kulturbevidstheden
med den „moderne Videnskab“ erklærer reentud: det
er umuligt! Det Overnaturlige er Skin og Phantasie.
Den Himmel, hvor den Almægtige throner, det Helvede,
hvori de Usalige pines, Lysets Engle, Afgrundens Dæmoner,
Altsammen. Skin og Phantasie!

Her er et Vendepunkt, afgjørende for Udviklingens
Gang og Slægternes Fremtid. Et Grundspørgsmaal er
% --- p. 11 ---
\setcounter{footnote}{0}%
\opage{11}tilbagetrængt, mishandlet, forkvaklet og kommer kun
frem i forvirrede Forestillinger; det er Spørgsmaalet om
Phantasiens Væsen og dens Betydning for Erkjendelsen
af den evige Sandhed.

Det Overnaturlige, hedder det, er kun en Phantasie.
Her er i denne almindelige Sætning Noget, vi maae
indrømme, og Noget, vi paa Liv og Død maae bekæmpe.
Vi maae indrømme, at alle de Forestillinger, vi kunne
gjøre os om et hiinsidigt, overnaturligt Rige, dannes i
Phantasiens Form, at de for saa vidt alle høre Phantasien
til. Men kun en Phantasie, dette lille Ord
kun er Stridens Kardinalpunkt, Kampens Midtpunkt,
dette kun bestemmer Seir og Nederlag, dette tilsyneladende
ubetydelige kun er som Waterlooslagets
„gamle Avlsgaard“.

Ved dette kun er der sat Splid mellem Virkelighed
og Phantasie. Her tænkes da udelukkende paa den
menneskelige Phantasie i dens mangfoldige deels uvilkaarlige,
deels vilkaarlige Formen, dens Gjæren og Væven,
Virken og Dannen. Det er da ganske sandt, at den
menneskelige Phantasie med al sin skabende Evne staaer
afmægtig overfor den virkelige Natur; den kan lege med
de naturlige Muligheder; den kan sidde ved Sandsynlighedernes
Væverstol og væve de usynlige Traade i
Mønstre, der have en skuffende, slaaende Lighed med
det Virkelige; men Virkeligheden selv kan den ikke
frembringe, den kan ikke skabe af Intet. Men hvad om
det ved nærmere Undersøgelse viste sig, at ikke blot
Tænkningens, men ogsaa Phantasiens Form er en oprindelig
og nødvendig Sandhedsform; at der ikke blot
gives en menneskelig, men ogsaa en guddommelig
Phantasie, ikke blot en afmægtig, men ogsaa en almægtig
% --- p. 12 ---
\setcounter{footnote}{0}%
\opage{12}Phantasie, ikke blot en eventyrlig spillende, men en i
Alvor villende, af Viisdommens Love, af Planer og Formaal
gjennemaandet Phantasie. Hvad om alt det, vi
kalde det Virkelige, det Haandgribelige, viste sig med
al sin Realitet at være, ikke et absolut Selvstændigt,
men et Billede præget af skabende, anskuende Tænken,
et Magtens Udslag af guddommelig Phantasie: det vilde
forandre Sagen. Kunde dette klares, begrundes og gjøres
indlysende for den dannede Bevidsthed, vilde da ikke
det Overnaturlige forlige sig med det Naturlige, og den
hele Tilværelse fremstille sig for os i et nyt Lys?

Her er ikke Tale om Brud paa Lovene; heller ikke
vil det for al menneskelig Viden uigjennemtrængelige
Mysterium blive gjennemsigtigt: Mystik og Mysticisme
maae vi holde borte. Men et nyt Mulighedernes Dyb
vil aabne sig for os. For at see ned i dette Dyb uden
at svimle, maae vi staae fast paa de strenge Loves
Grundlag. Det er en ny Himmel og en ny Jord, vi
skulle opdage, ikke i Virkelighedens Form, men i Mulighedens.
Da Alt skal sees under Mulighedens Slør i
Phantasiens Speil, er det nødvendigt at kjende den
Grændselinie, henad hvilken vi skulle bevæge os: at
klare Phantasiens Væsen er ikke det Samme som at
phantasere. Vi maae have Virkelighedens Blylod med,
for at prøve, hvor dybe Grunde der er i Mulighedernes
Ocean. Det turde da vise sig, at der kan være Verdener
af anden og høiere Orden end dem, der tilhøre
dette naturlige Universum.

Vi kunne ikke følge Beatrices Elsker paa hans
høie Flugt; vi maae gaae Undersøgelsens jevne Vei og
bruge vore Maalesnore; vi maae skjelne mellem Illusion
og Virkelighed, mellem Skin og Sandhed. Poesien har
% --- p. 13 ---
\setcounter{footnote}{0}%
\opage{13}sin Phantasie, men Religionen har ogsaa sin. Phantasien
behersker Menneskene lige saa vel middelbart
som umiddelbart. Vi ville forsøge en Vandring gjennem
det store Rige, hvor den Vidunderlige boer og aabenbarer
sin Magt; dette Rige er ikke langt borte; tvertimod,
det er os uendelig nær. Mange af Phantasiens
Phænomener ere os bekjendte; men Væsenet er et
Ubekjendt: vi ville stræbe at gjøre dette Ubekjendtes
nærmere Bekjendtskab.

% --- p. 14 ---
\setcounter{footnote}{0}%
\lecture{Anden Forelæsning.}\opage{14}

Der er i vor Tilværelse neppe noget Emne af høiere
Rang, der i den Grad har formaaet at hendrage det
store Fleertals Opmærksomhed paa sig som Spørgsmaalet
om Sjælens Udødelighed. Det er et vedvarende
Stridsspørgsmaal. Striden føres ikke alene udvortes
imellem dem, der afgjørende bekræfte, og dem, der lige
saa afgjørende benegte Udødeligheden; men Stridsemnet
lister sig ind i hvert Menneskes Indre; det
vækker en Splid mellem Sands og Forstand, mellem
Hoved og Hjerte, mellem Tro og Tvivl. Det er et vedholdende
Bekymringsspørgsmaal; det suger sig fast i
Bevidstheden, det trænger gjennem Blod og Nerver ind
i Samvittigheden. Hvor sandselig sløvt eller hvor
tankeløst og letsindigt maa ikke det Menneske være,
hvem dette Spørgsmaal om Sjælen og dens Udødelighed
aldrig berørte. Det er mere end et almindeligt Tankespørgsmal, % PRINTED AS IS: „Tanke-|spørgsmal“ (an a missing)
det er personligt. Fra dette Stridsemne er
det umuligt at udelukke Phantasien; den blander sig i
Sagen, enten man mærker det eller ikke; ja den
kommer tilsidst uvilkaarlig til at spille en Hovedrolle.
Man lægge Texten tilrette, hvordan man vil, det er
dog Phantasien, der skriver Nodern og sætter Musikken % PRINTED AS IS: „Nodern“ (a gap: the final e did not print)
% --- p. 15 ---
\setcounter{footnote}{0}%
\opage{15}til. Intet Thema har flere Variationer. Hvor bliver
der ikke phantaseret over Sjælens Udødelighed!

Vistnok falder Spørgsmaalet om Sjælens Udødelighed
oprindelig sammen med Spørgsmaalet om det evige
Liv, og bliver for saa vidt religiøst at besvare i Tro;
men ogsaa Troen har sin Phantasie. Imellem Tro og
Phantasie er der et eget Forhold. Troen vil lige saa
lidt som Phantasien underordne sig det udvortes Sandselige;
begge ere enige om, ikke at krybe for Naturlovene
eller bøie Nakken under den blinde Nødvendighed.
Begge føle sig fribaarne, i Stand til at optage Kampen
med det Endelige og hæve sig over de Skranker, der
sættes dem af den materielle Virkelighed. Og dog kan
Troen ikke uden Betænkelighed give sig i Ledtog med
Phantasien. Den aner en Fare. Der er i Phantasien
en Drift til at gaae paa Eventyr, en ungdommelig
Kaadhed, en Lyst til at drille Virkeligheden og lege
Sandsynlighedsleg med de store Muligheder. Paa denne
Leg kan Troen ikke indlade sig, thi dens Væsen er
Alvor; Troen kan, om saa skal være, byde den timelige
Verden Trods, thi den kjender dens Forkrænkelighed
og har gjennemskuet dens Intethed; men Virkeligheden
selv kan den ikke byde Trods; tvertimod! den er
uendelig interesseret i det Virkelige, det høiere Virkelige,
for hvilket den strider. Troen tør ikke kaste
Virkelighedens Byrder fra sig, for at gjøre sig let og
svinge sig høit paa Phantasiens Vinger; thi det er
Ikarusvinger, Vinger af Vox, der smelte for Solens Ild.

Men, seet fra en anden Side, ville Tro og Phantasie
gjøre Forbund, nemlig hvor det gjælder at komme til
Rette med det Sandsynlige og erkjende det Mulige.
Troen kan gaae imod Sandsynligheden; men den kan
% --- p. 16 ---
\setcounter{footnote}{0}%
\opage{16}ikke gaae imod Muligheden. Det Usandsynlige kan
troes, men det Umulige kan ikke troes; saa maa de
begge, Troen og Phantasien, hjælpe hinanden med at
udklare Muligheden.

Anvendes det her Sagte paa Læren om Sjælens Udødelighed,
da er det Første og Vigtigste dog vel dette, om en
saadan Udødelighed er mulig. Dette Spørgsmaal bliver
bestandig mere nærgaaende, alt som vi lære denne
virkelige Verden, den nærværende Natur at kjende.
For Oldtiden var Spørgsmaalet ikke nær saa vanskeligt
eller indviklet som det er for Nutiden; thi Nutiden har
en anderledes indtrængende Viden om Naturlovene og
det Materielles Betydning end Oldtiden. Der er i vor
Tid Mennesker, som leve i hemmelig Angst for Naturlærens
voxende Indflydelse; de have en Frygt for at
indlade sig med dens Undersøgelser, de ere bange for
at miste deres Tro. Det er en piinlig Tilstand og et
sørgeligt Vidnesbyrd, at det nok ikke staaer ganske
rigtig til med denne Tro, at der i Hjertet er en Tvivlens
Orm, som gnaver. Har Troen Sandhed, behøver
den da at leve i et Tusmørke? den maa kunne taale
Dagens Lys. Der er slet ingen Fare ved at lade Naturlæren
rykke frem med sine Indsigelser; det skal nok
vise sig, at den selv er i Vildfarelse og gjør Overgreb,
naar den bilder sig ind at kunne udslette Muligheden,
den Mulighed, hvori Troen grunder. Synes det end et
Øieblik at lykkes, da er det kun i Kraft af en ufuldbaaren
Tænkning og en misledet Phantasie.

Vi ville anføre nogle af de bekjendteste Indvendinger
imod Sjælens Udødelighed. Det er umuligt, siger
man, at Sjælen kan overleve sit Legeme, thi det sjælelige
Liv er afhængigt af Organer. Det vi kalde vor
% --- p. 17 ---
\setcounter{footnote}{0}%
\opage{17}Sjæl er ikke noget Væsen for sig. Vel er Sjælen ikke
Materie; men den er dog en Virksomhed af det Materielle,
en Nervevirksomhed og fremfor Alt en Hjernevirksomhed.
Hvor bliver Sjælen af, naar den forlader
Legemet? Ja, hvor bliver Lyset af, naar Ætheren
kommer i Hvile; hvor bliver den elektriske Strøm af,
naar Batteriet forstyrres; hvor bliver Ilden af, naar
den chemiske Proces hører op? Det forholder sig med
Sjæl og Legeme som med Tonerne og et Instrument.
Et forstemt Instrument giver falske Toner, og naar
Strengene briste, er det forbi med Musikken. Det er
en barnagtig Forestilling, at Sjælen skulde forlade Legemet
eller, som det hedder, fare ud af Legemet. En
nøgen, legemløs Sjæl er en tom Phantasie. I Orienten
troer man paa Sjælevandringen; de Afdødes Sjæle kunne
boe i Planter og Dyr, i Katte, Hunde, Elephanter, ja
den troende Hindu maa endog være bange for at kunne
træffe sine Forfædres Aander i Rotter og Muus: ængstende
Indbildninger, uhyggelige Phantasier. Hos os
hører Troen paa Sjælevandringen vel til Sjeldenhederne,
men desto almindeligere er Forestillingen om, at Sjælen
gaaer over til en anden Verden. Hvor ligger nu den
Verden? Phantasien kommer i Forlegenhed, naar den
examineres af Forstanden. Den digter en hiinsidig
Verden, men hvor er dette Hiinsidige? Bagved, udenfor,
indenfor, nedenfor, ovenfor den nærværende Verden?
Disse Spørgsmaal kan Phantasien ikke besvare. Hvorsomhelst
den peger hen, opad eller nedad, i Nord eller
Syd, i Øst eller Vest, den peger allevegne ud i Rummet
og kommer ikke udenfor dette naturlige Universum.
Der phantaseres om et nyt Legeme, et aandeligt, overnaturligt
Legeme, altsaa en aandelig, overnaturlig Ma%
% --- p. 18 ---
\setcounter{footnote}{0}%
\opage{18}terie. Men er der noget Naturligt til, saa er det nok
Materien. Hvad er vel en overnaturlig Materie? En
unaturlig Natur, et udstrakt Væsen uden Udstrækning.
Man har desaarsag i den nyere Tid da ogsaa begyndt
at indrette „Dagen efter Døden“ paa en med Naturlæren
lidt mere overeensstemmende Maade ved at betjene
sig af Brinten og Ætheren og lade Sjælen afreise
paa en Lysstraale til andre Kloder, Syvstjernen eller
Nordstjernen eller maaskee til Hundestjernen. Det
smager af Videnskab; men Smagen er vammel, og der
behøves kun ringe Eftertanke for at overbevise sig om,
at det Hele er et løsagtigt Phantasterie.

Naar disse Indvendinger udvikles og gjennemføres
med skarp Forstand og Kjendskab til Naturlovene,
gjøre de en stærk Virkning. Jo mere Overveielsen gaaer
i det Enkelte, desto urimeligere bliver det Hele. Det
Sande og det Falske blandes paa en saa eiendommelig
Maade, at Sagen let kan faae et Udseende, som var
Udødelighedstroen kun en Phantasie, en Indbildning.
Hvoraf kommer nu det? Det kommer deraf, at Grundlaget,
hvorpaa den hele Betragtning hviler, er falsk.

Det er for det Første en falsk, en grundfalsk Forestilling,
at Sjælen ikke er et fra Legemet forskjelligt
Væsen. For at bevise, at Følelse, Bevidsthed, Forstand,
Phantasie ikke er Andet end Nerve- og Hjernevirksomhed,
gaaer man ud fra Erfaring; men Erfaringen kommer i
Strid med sig selv; thi der gives ikke blot en udvortes,
men ogsaa en indvortes Erfaring. Jeg erfarer, at jeg
har et Legeme, og at dette Legeme bestaaer af materielle
Dele; men jeg erfarer ogsaa, at der er Noget i mit
Indre, jeg maa kalde mig selv; jeg erfarer, at min Bevidsthed,
min Følelse, min Forstand og min Phantasie,
% --- p. 19 ---
\setcounter{footnote}{0}%
\opage{19}hvormeget disse Virksomheder end ere betingede og
afhængige af materielle Organer, af Nerver og Hjerne,
dog umulig kunne være materielle Virksomheder. Magneten
kan tiltrække, et elektrisk Batteri kan udvikle
Elektricitet, Sligt hører til de livløse Elementers Natur.
Men Materien, hvor eiendommelig den end fordeles og
leires, kan ikke føle; der tænkes ved Hjælp af Hjernen,
men Hjernen, denne besynderlig furede Masse med sine
Vindinger og Labyrinther, Hjernen som blot organiseret
Materie kan ikke tænke. Det er vor Sjæl, der føler,
det er vor Sjæl, der er sig bevidst, det er vort sjælelige
Væsen, der tænker. Kan det nu nytte at komme
med den Forklaring, at det med Sjælen kun er en
Phantasie. Hvorfra kommer da Phantasien? Vi gad
høre den Physiker, som kunde forklare os med nogen
Rimelighed, hvorledes Atomerne maae leire sig, for at
Materien kan begynde at phantasere. Man forsøge det
en Dag eller endog blot en Time at leve med Materialisternes
eller de moderne Positivisters Anskuelse, og man
vil gradviis nærme sig Vanvid. Skal Sjælen være Hjerne
og Hjernen Sjæl, saa maa den hele menneskelige Organisme
betragtes under Synspunktet af en eiendommelig
formet Organdynge; man kan da ikke tale om mit eller
dit Legeme, om min eller din Hjerne, thi Udtryk som
min og din forudsætte, at Du og jeg ere særlige Væsener.
Naar to Mennesker tale med hinanden, er det altsaa
ikke to Sjæle, men to Hjerner, der meddele hinanden
deres Forestillinger og betroe hinanden deres Tanker.
Det kan jo for et ironisk Lune være ret underholdende,
at betragte Menneskene paa Torve og Gader som selvbevægende
Atomdynger, organiske Automater, at see
dem nikke til hinanden og tage Hatten af for hinanden,
% --- p. 20 ---
\setcounter{footnote}{0}%
\opage{20}idet enhver af disse Hjerner bilder sig ind, at den har
et Selv, et Jeg, skjøndt alt det med Jeg selv og
Du selv kun er Phantasie. I en af Hoffmanns Romaner
skildres en ung Pige, som ved sin Skjønhed
og sine Talenter gjør overordenlig Lykke i det elegante
Selskab; hun dandser graciøst, hun synger og
spiller med Virtuositet, med Smag og Foredrag. De
unge Herrer flokkes om hende, hun fængsler dem i en
Tryllekreds; der er Tilbedelse og Courtoisie, der er
Hjertebanken og Concurrence; der er lynende Blikke,
som endogsaa true med Dueller. Men En af Tilbederne
holder sig stille tilbage, han venter. Han venter paa
det Øieblik, da han kan træffe den Skjønne alene og
kaste sig for hendes Fødder. Øieblikket kommer, men
hvad maa han erfare? Den Tilbedte er intet virkeligt
Menneske; det er en Dukke, et overordenlig fiint forarbeidet
Konstværk! Øinene af slebet Kul, de mørke
Bryn paasatte, hun er en Automat, hun har ingen Sjæl,
hun gaaer paa Skruer. Hvad om vi nu alle vare naturfrembragte
Automater, der bildte os selv og Andre ind,
at vi havde en levende Sjæl, medens dog Sandheden
var, at vi alle gik paa Skruer. Vanvid!

Det maa staae fast for den sunde Sands og den
indre Erfaring, at vi have en Sjæl, og at den i vort
Legeme virksomme Sjæl er et fra Materien forskjelligt
Væsen; og saa er det da ikke, hvad Materialisterne
ansee det for, barnagtigt at tale om, at Sjælen ved
Døden forlader sit Legeme. Men hvor gaaer den hen?
Her kommer Phantasien til at tale et Ord med, og det
for Alvor. Hvad er Rummet, hvad er Materien, hvad
er den legemlige Verden? Uden Phantasiens Hjælp vil
% --- p. 21 ---
\setcounter{footnote}{0}%
\opage{21}det være umuligt at give noget forstaaeligt, end sige
noget fornuftigt, Svar paa disse Spørgsmaal.

Hvad er Rummet? en Form for det Udstrakte, i
sig selv en tom Form, men dog en Anskuelsens Form,
en Form for alt det Anskuelige. Dette Anskuelige kan
nu være deels et Udvortes og Materielt, deels et Indvortes
og Aandeligt; men Rummet retter sig altid efter
det Indhold, hvormed det fyldes. Den, der drømmer,
drømmer i Rummet; men det er ikke det udvortes Rum,
det er et indvortes, og dog faaer dette Indvortes Form
af et Udvortes. Den Bortreiste f. Ex., der længes, kan
ligge paa et fremmed Sted i det fjerne, fremmede Land
og drømme om sit Hjem. Drømmen kan være saa
levende, at han seer dem alle for sig i de bekjendte
Værelser og de vel bekjendte Omgivelser. Der er altsaa
et andet Rum i hans Drøm end det, hvori han
legemlig er, og hvori han vil befinde sig, naar han
vaagner. Her er altsaa to Slags Rum, det drømte og
det virkelige. Hvorfor kalde vi nu ikke det drømte
Rum virkeligt? Fordi det ikke er fyldt med virkelige
Gjenstande, men kun med Drømmebilleder. Som den
tomme Form er det drømte Rum lige saa vel Rum som
det virkelige: Forskjellen beroer alene paa det, hvormed
Rummet fyldes. Et Maleri paa Væggen, et Landskab
f. Ex., viser os det Samme paa en anden Maade.
Rummet i Værelset og Rummet i Billedet give to Slags
Rum. Billedet, det forestille nu den romerske Campagne
eller den jydske Hede, har sit eget Rum. Forgrund,
Mellemgrund, Baggrund skille sig ad; det Nærmere
nærmer sig, det Fjernere fjerner sig: vi see ind i
en ubestemmelig Dybde, og dog bliver Væggen ikke
gjennembrudt. Hvad er nu det for et Rum? Et illu%
% --- p. 22 ---
\setcounter{footnote}{0}%
\opage{22}sorisk Rum og hvorfor? Fordi de Gjenstande, hvormed
det fyldes, ere illusoriske.

Phantasien er altsaa ikke i Modsigelse med sig
selv, naar der skjelnes imellem forskjellige Slags Rum.
Hele Vanskeligheden ved Her og Hisset falder bort,
hvad Rummet angaaer. Vi maae da udelukkende spørge
om det, der opfylder Rummet. Vi komme herved til
at gjøre et dybere Indblik i Phantasiens Væsen.

Vor menneskelige Phantasie kan fylde et indvortes
Rum med Billeder, men Rummet forbliver illusorisk,
et indbildt Rum, fordi vore Billeder ere Indbildninger,
der mangle Virkelighed; det er vor Phantasies Afmagt,
at den ikke kan give sine Billeder sand Realitet. Hvad
hører der altsaa til for at frembringe et virkeligt Rum?
En almægtig Phantasie, der kan fylde sit Rum med
virkelige Gjenstande, en skabende Phantasie, der formaaer
at oversætte sit rige Indhold, sit Indre, i et
mangfoldigt og mægtigt Ydre, d. e. i en udvortes, sandselig
tilværende Verden.

Det er Forholdet imellem det Indre og det Ydre,
hvorpaa det her kommer an. Vi tale om Skaberens
almægtige Tanker og om hans almægtige Villie; men
lad os ikke glemme hans almægtige Syn. Dersom der
i Skaberens evige Selvmagt ikke var Andet end Tanker
og Villie, saa var denne udvortes, legemlige Verden
aldrig blevet til. Vel er der Tanker, guddommelige
Tanker i denne Verden, men her er ogsaa Materie og
Legemlighed; denne Materiens og Legemlighedens Verden
er mere end en Verden af Tanker. Tanker ere usynlige,
men her er en Verden af haandgribelige Ting.
Vel er der i denne Verdens Love en lovgivende Villie
kjendelig; men der er i Verden mere end Love, der er
% --- p. 23 ---
\setcounter{footnote}{0}%
\opage{23}Stoffer og Masser. Lovene ere usynlige som Villien,
hvoraf de fremgaae; men den Tingenes Orden, hvori
Lovene opfyldes og Villien fuldbyrdes, er sandselig,
synlig, haandgribelig. Hvorfra stammer nu dette Udvortes,
Sandselige, Haandgribelige? Fra Phantasien, den
almægtige Phantasie. Det er jo netop Phantasiens
Væsen at være mæglende imellem det Indre og det
Ydre, det Usynlige og det Synlige. Phantasien er beslægtet
med Tanken og tilhører det Indre; men den er
ogsaa beslægtet med det Sandselige, særlig med Synets
Sands: dens Frembringelser ere Billeder, og naar der i
disse Billeder er Magt, saa ere Billederne virkelige
Gjenstande. Magt og Virkelighed, Magt og Selvstændighed,
Magt og Realitet svare til hinanden i det
Mindste som i det Største. Der er en Selvstændighed
i Tingene, fordi der er en Magt i Tingene; der er en
Selvstændighed i hvert Atom, fordi der er Noget af en
Magt i Atomet: jo mere Magt, desto større Selvstændighed,
jo mindre Magt, desto mindre Selvstændighed.
Den hele Natur er et Magtens Rige. Derfor har Naturen
en Sluttethed i sig selv, den udvikler sig af sit
Indre, sine Anlæg; men alle Bestemmelser ere forud
givne, den kan ikke give sig selv et eneste nyt Anlæg.
Denne Magt i Naturen stammer fra Almagten; den er
en Synliggjørelse i guddommelig Phantasie, en Aabenbarelse
af den uendelige, skabende og opholdende
Selvmagt.

Hvor mange Verdener kan der altsaa være? Vi
svare: saa mange som der kan rummes i den guddommelige
Phantasie, og det er --- mange!

Frelserens Ord: „I min Faders Huus ere mange
Boliger“ tyde hen paa, at hans Blik, om det end
% --- p. 24 ---
\setcounter{footnote}{0}%
\opage{24}maaskee har hvilet paa den østerlandske Stjernehimmel,
paa Pleiaderne eller Orions Bælte, dog har gjennemtrængt
dette Ydre og seet ind i en anden Himmel, en Verden
af høiere Orden. Der kan altsaa foruden dette vort
synlige Universum med sin Uendelighed, sit Rum, sin
Materie, sin Legemlighed og sine talløse Kloder, være
utallige høiere Verdener, hver især med sin Uendelighed,
sit Rum og sin Legemlighed. Vor Stjernehimmel er
høi og dyb: men dybere endnu er den almægtige Phantasie
med de uudtømmelige Muligheder. Fra dette
Synspunkt ville vi betragte Sjælens Udødelighed.

Hvad vil det nu sige, at Sjælen forlader Legemet?
Vi udtrykke det tydeligst ved, at den ombytter den
ene Verden med den anden. For at indsee, hvad heri
ligger, kaste vi et Blik paa Søvnen og Døden. Døden,
hedder det, er Søvnens Broder; Søvnens Genius med
den opadvendte Fakkel svarer til Dødens med den
nedadvendte. Dette Broderskab gjælder ikke i physiologisk
Forstand; Physikeren og Lægen ville have Meget
at indvende derimod; men det gjælder i psychologisk
Forstand; det gjælder, naar der spørges om det Indre,
om Livets Væsen og Grund. I Søvnen lukke sig de
udvortes Sandser, Bevidstheden tager af, og Søvnen kan
være saa dyb, at endogsaa Drømmen forsvinder. I
Døden ophører al Sandsning, Bevidsthedens Lys slukkes;
men Sjælen bliver ikke til Intet; den synker i sit Indre
ned i sin dybe Grund; denne Grund er det evige Væsen,
den guddommelige Phantasie, i hvis Skjød den hviler.
Derfor nævne vi de Døde med det betegnende Ord, de
Hensovede. Paa Dødssøvnen følger en Opvaagnen, en
ny Legemdannelse i en ny Verden. Den nye Verden
er af anden og høiere Orden end denne nærværende;
% --- p. 25 ---
\setcounter{footnote}{0}%
\opage{25}men baade den Verden, som Sjælen forlader, og den,
hvori den ved Dannelsen af sit nye Legeme indtræder,
begge Verdener, hvor forskjelligartede de end ere, ere
Aabenbarelser af den ene og samme almægtige Phantasie.
Med Opvaagnelsen følger Erindringen om det
Forbigangne; men Overgangen er ingen Reise, den er
en Forvandling. Overgangen er indvortes, og dog virkelig;
vi kalde vore Døde de Hedengangne.

Saa langt kunne vi gaae uden at støde an mod
nogensomhelst af de bekjendte Naturlove; men vove
vi os videre frem, forsøge vi at skildre, hvorledes det
seer ud i den nye Verden med de nye Tilstande, saa
begynde vi at phantasere.

Der er hos visse Mennesker en ubændig nysgjerrig
Higen efter at gjøre fortroligt Bekjendtskab med Aandeverdenen,
især naar det kan skee paa en ret sandselig
og haandgribelig Maade; de mane Aander frem, de
adspørge de Døde. Det gik vidt i Oldtiden, man tænke
paa „Keiser og Galilæer“; men det gaaer endnu langt
videre i Nutiden. Der er i denne Henseende intet
Forehavende saa dumdristigt, at en ophidset Phantasie
jo tør indlade sig derpaa. At faae Aanderne til at
vise sig legemlig, især i Mørke, er nu som Ingenting;
men at udfritte dem ret privat, at photographere dem,
at tage dem med i Huusholdningen og faae dem til at
lave Pølser, det er --- „en Videnskab“. Oprørende er
det, naar et saadant Phantasterie af den groveste Aandsraahed
vil sætte sig i Forbindelse med Religionen og
hjælpe til med at faae Christendommen videnskabelig
begrundet.

Der er en Form af Phantasie, som staaer Religionen
nær, den er af ædlere, høiere Oprindelse, den har sit
% --- p. 26 ---
\setcounter{footnote}{0}%
\opage{26}Udspring, ikke fra den grove Sandselighed, men fra en
frit tilegnet, klar Anskuelse af Livets Kriser; det er en
Digterphantasie, der udfolder sit Indhold i den saakaldte
apokalyptiske Poesi. Herhen hører som Led i
Verdensliteraturen Dantes divina Comedia, i vor egen
Literatur Heibergs „En Sjæl efter Døden“; ogsaa Slutningspartiet
af „Adam Homo“ bør nævnes. De Billeder,
som forekomme i slige Digtninger, gjøre ikke Fordring
paa at gjælde for Virkelighed, det er Symboler, De
kunne være af gribende Virkning, og faae ethisk Betydning
derved, at de anskueliggjøre Sjæletilstande, som
paa en indtrængende Maade minde os om de personlige
Udviklingslove og levendegjøre vor Forestilling om
Sammenhængen mellem Livet her og Livet hisset. Man
tænke paa Adam Homo, som han sidder paa Synderbænken og seer med Gru paa den store Vægtskaal.
Alma nærmer sig, en Haabets Stjerne i den mørke Nat;
hun trøster den Forsagte, leder ham ved Haanden og
indgyder ham Mod til at gjennemgaae Ildprøven. Det
er en Strid mellem det Nye og det Gamle, en Forvandlingsproces,
som gjøres anskuelig i Situation. Den
forhen Store bliver saa lille, og den forhen Lille saa
stor. Adam Homo, i Verden den store Mand, „Adam
Homo med den blide Aand, Baron, Geheimeraad, Ridder
af det hvide Baand“, hvor er han blevet lille; og Alma,
den stille, dybe, men oversete og forladte Sjæl, hvor er
hun stor! For at den nye Herlighed kan aabenbares,
er der et Jordstof, der først maa fortæres, et „Pandser“,
der maa smelte, en gammel Herlighed, der maa
„brændes til Kul“. Og det foregaaer for vore Øjne,
det bliver os nærværende, det lyder ud fra Skjærsildens
Chor: „Hjertets smeltede Pandser henflyder, Med
% --- p. 27 ---
\setcounter{footnote}{0}%
\opage{27}vor Pynt drive Flammerne Spøg: Vore Krandse, Bedrifter
og Dyder gaae op som et Intet i Røg“. Det er
Phantasie, men denne Phantasie er under Ideens Tugt,
det er ingen ørkesløs, tom Phantaseren.

Og dog er der en Forskjel, en væsenlig Forskjel
mellem den apokalyptiske Poesi og Aabenbaringsordet.
Aabenbaringsordet anskueliggjør i store Træk, hvad der
venter os i en tilkommende Verden. Ogsaa disse Skildringer
høre ind under  Phantasiens Omraade; men her
maa anlægges en ganske anden Maalestok: her sigtes
i høiere Forstand paa et Virkeligt. Her er intet Grundlag
for videre gaaende Phantaseren, men en Opfordring
til alvorlig Selvprøvelse. Vi skulle i det Følgende gaae
over til at betragte den religiøse Phantasie.

% --- p. 28 ---
\setcounter{footnote}{0}%
\lecture{Tredie Forelæsning.}\opage{28}

Alt for en Deel Aar tilbage blev der i vor Literatur
gjort opmærksom paa, at Religionen led under Mishandling
af Videnskaben. Det blev da fremhævet, at
der maatte gjøres Forskjel mellem Naturligt og Overnaturligt,
at Aabenbarelser fra et hiinsidigt Rige ikke
lode sig maale med Hverdagserfaringers Alen eller
presse ind under de bekjendte Naturlove. Det blev
fremhævet, at ihvorvel Intet kan være sandt i Religionen,
hvis Usandhed videnskabelig lader sig bevise, saa var
det dog utilbørligt at sætte Videnskaben til Dommer
over Religionen; det var at forgude vor stykkevise
Viden og løfte den op paa Alvidenhedens Throne.
Religionen maatte have Lov til at udvikle og begrunde
sine Sandheder religiøst; saa kunde Videnskaben paasee,
at der i Udviklingen Intet indeholdtes, som stred mod
vitterlige Kjendsgjerninger og i alt Fald varetage sin
ubestridelige Ret til alsidig fri Undersøgelse. Det blev
i udførlig Fremstilling viist, at der ifølge den menneskelige
Videns Begrændsning nødvendig maatte gives to,
af hinanden uafhængige, men dog hinanden vexelsidig
belysende Erkjendelseskredse, en for det Overnaturlige
og en for det Naturlige, en for Religionen og en for
% --- p. 29 ---
\setcounter{footnote}{0}%
\opage{29}Videnskaben. Der blev, kort sagt, gjort en utvetydig
Adskillelse mellem Tro og Viden.

Denne Adskillelse vakte Forargelse, ikke blot hos
dem, der paa Forhaand holdt sig forvissede om, at
Troens Dage vare talte, eftersom det Overnaturlige,
betragtet i Videns Lys, kun er Skin og Phantasie; men
hvad der er mere mærkværdigt, ogsaa hos dem, som
paa Kirkelærens Vegne følte sig kaldede til at tale
Religionens Sag. De erklærede reent ud, at det Overnaturlige
ikke alene lod sig optage til Begrundelse i
Videnskaben, men at der gaves en ifølge sin Oprindelse
overnaturlig Videnskab, en Hellig-Aands Videnskab, der
satte sine Dyrkere i Stand til at hæve sig over „Verdens
Børnelærdomme“ og tugte, „de hedensk opblæste
Naturvidenskaber“. Og hvad Vei gik man saa? Ikke
Logikens Vei, af Frygt for at smitte sig med „det
Kenologiske“; man valgte en mere anskuelig Fremgangsmaade,
nemlig at lade saadanne Ord som Villie,
Skabervillie, den frie Inderlighed, den gode Villie, aftrykke
med fede Typer. Fede Typer ere jo ogsaa
Realiteter; og ved slige Realiteter forberedte man sig
da til at tugte de indbildske Realvidenskaber. Som
sagt, Forargelsen over, at Religionen skulde have sin
egen Erkjendelseskreds var paa begge Sider saa stor,
at Troens Forsvarere og Troens Fornegtere, Kirkens
Mænd og Mænd af den frie Tanke i Videnskabens Navn
rakte hinanden Haand, og den gjensidige Forstaaelse
syntes at blive cordial. Det Cordiale tabte sig dog, og
det Hele gik lidt efter lidt over til en stille Gjæring.
Den ærede Almeenhed har siden den Tid været bekymret
snart for Videnskaben og snart for Religionen,
snart for „Troen“ og snart for „den frie Tanke“; men
% --- p. 30 ---
\setcounter{footnote}{0}%
\opage{30}i de sidste Dage ere Udsigterne blevne lysere; rette
Vedkommende, som tilgavns ere inde i Sagen, have
omsider paataget sig det besværlige Hverv at oplyse
det store Publikum. Den bekymrede Almeenhed vil
nu kunne sætte sig i Kundskab om, hvorledes det
egenlig staaer til med det omstridte Mellemværende
mellem Religion og Videnskab. At en Afgjørelse nærmer
sig, er øiensynligt, det kan tydelig sees deraf, at Striden
føres i Aviserne. Naar en saadan Strid om en saadan
Sag gaaer over i Aviserne, saa veed man, at nu er det
Alvor; det er et ligesaa ubedrageligt Tegn paa Alvor,
at en Sag af denne Natur kommer i Avisen, som det
er et ubedrageligt Tegn paa, at en Melodie er almeenyndet,
naar den kommer i Lirekassen.

Vi kunne ikke her tage hele Sagen for i den
Korthed, den Fynd og Klarhed, hvormed den bør tages
i Avisen; Sligt overlade vi til dem, som dertil have
Begavelse, Forstudier og Myndighed. Men et Bidrag
kunne vi dog give, nemlig til at belyse den Side, som
særlig angaaer Phantasien.

Den aabenbarede Religion er grundet paa Mirakler;
men, derfor indestaaer „den moderne Videnskab“, Miraklet
er kun Skin og Phantasie. Der skeer ingen
Mirakler og der er aldrig skeet Mirakler; Miraklet er
en selvmodsigende og desaarsag intetsigende Forestilling.
Miraklet er et Brud paa Naturlovene, men Naturlovene
ere uforanderlige; det er en Ramse, som enhver Skoledreng
nuomstunder kan paa sine Fingre. --- Men hvad
er nu egenlig Meningen med, at Lovene brydes? Meningen
er, at Tingenes Sammenhæng brydes, at den
guddommelige Viisdom modsiger sig selv. Denne Sætning
har Gyldighed paa een Betingelse, den nemlig, at
% --- p. 31 ---
\setcounter{footnote}{0}%
\opage{31}hele Viisdommens Plan er udtømt, fuldstændig udtømt
i denne naturlige Tingenes Orden. Kan Videnskaben
virkelig bevise dette --- endnu har den ikke beviist det;
dens foregivne Beviser ere kun Paastande, og disse
Paastande i maadelig Overeensstemmelse med dens
ædruelige Bekjendelser, at den ikke kjender Væsenet,
men kun Egenskaberne, ikke Kraften selv, men kun
dens Yttringer o. s. v. --- men skulde Videnskaben engang
naae til en saa dybtgaaende Indsigt i Tingenes
Væsen og hemmelige Sammenhæng, at den med klare
Grunde kan godtgjøre, at den overnaturlige Verden
kun er Skin og Phantasie; saa er det i Sandhed forbi
med Aabenbaring og Tro. Men heller ikke før, lad os
ikke forhaste os!

Aabenbaringsfornegterne ere snilde nok til at holde
Striden fast paa Naturlovenes Omraade og lade, som
om det med det Overnaturlige for længst var afgjort:
paa denne Maade have de allerede Seiren i Hænde.
Hvad kan det hjælpe, at Troesforsvarerne beraabe sig
paa den almægtige Skaber, naar Skaberens Plan med
Love og Sammenhæng, Materie og Kræfter er udtømt,
fuldstændig udtømt i dette naturlige Universum? Er
der Intet bagved, slet Intet, ingen fri Inderlighed, intet
Rigdoms Dyb af Muligheder, hvoraf den skabende
Frihed kan øse, hvad nytter det saa at paakalde Almagten?
Den evige Lov er Lovgiverens Herre, Nødvendigheden
raader: Frihedens Puls staaer stille. Den
menneskelige Villie har dog en Art fri Inderlighed, en
Selvbestemmen; men er Guddommens Skjød umiddelbart
Eet med Naturens Grund, da bryder Alting frem,
ja Livet selv, af betinget Nødvendighed, som naar en
Gnist falder i Krudt. Vi leve i Constitutionernes
% --- p. 32 ---
\setcounter{footnote}{0}%
\opage{32}Tidsalder, Videnskaben har faaet Øiet op for Constitutionen
i det store Hele; den vil Intet høre om „en
uansvarlig Gud“ med urandsagelige Raadslutninger; den
har gjennemseet Tilværelsens gamle Grundlov paany
og forvisset sig om, at den Almægtige nok skal afholde
sig fra at gjøre Statscoup.

Lad det være en Phantasie, saa staaer det Spørgsmaal
dog tilbage, om de kirkelig beskikkede Troesforsvarere
fra deres Side virkelig formaae at indføre
det Overnaturlige i Videnskaben og faae dets Erkjendelse
tilstrækkelig begrundet. At tage Myndighedsminen
paa og bruge Magtsprog, hjælper ikke; at ville paatvinge
Videnskabernes frie Republik et absolut Monarchie
og lade „den gode Villie“ trykke op med nok saa
fede Typer, for at tugte de naturalistiske Børnelærdomme,
fører til ydmygende Nederlag. Her er en Phantasiens
Hemmelighed, vi maae see at faae opklaret.

Den aabenbarede Religion har Syner og Billeder,
underfulde Former af skabende Phantasie; vel at mærke
ikke blot af menneskelig, nei, oprindelig af guddommelig
Phantasie. Den lange, tusindaarige Strid om Religion
og Videnskab bærer Vidne herom. De troende Fædre
beraabte sig paa Aandens Vidnesbyrd, og med Rette.
Men de opfattede Vidnesbyrdet eensidig, idet de ikke
nok som agtede paa den religiøse Phantasie, og desaarsag
overvurderede det udvortes, historisk Virkelige.
Det gik til en Tid. Men saa kom der et Vendepunkt.
Alt som Forstandsoplysningen tiltog, og Kritiken begyndte
at komme til Kræfter, indtraadte Vanskeligheder,
den ene mere umedgjørlig end den anden. De
protestantisk-orthodoxe Lærere, der i Modsætning til
Katholikerne meente at have et tilforladeligt juridisk
% --- p. 33 ---
\setcounter{footnote}{0}%
\opage{33}thinglyst Aabenbaringsdokument i den hellige Skrift,
gik ud fra Forudsætningen om Skriftens absolute Ufeilbarhed.
For at forherlige Ufeilbarheden antoge de, at
hvert Ord, hvert Bogstav og Tøddel saavel i det Gamle
som i det Nye Testamente var indblæst af Gud; denne
Antagelse beredte dem smertelige Skuffelser. Til Ufeilbarheden
hørte, at ogsaa Vocaltegnene i den hebraiske
Text oprindelig maatte være indblæste; men Kritiken
viste, at disse Tegn vare en sildigere Tilsætning af
jødiske Lærde, og Orthodoxiens Ordførere maatte give
efter. Til Ufeilbarheden hørte, at Græciteten i det Nye
Testamente maatte være classisk, thi den Aand, der
paa overnaturlig Maade havde uddeelt sine Sproggaver
paa Pintsefesten, maatte dog vel være lige saa grundig
hjemme i det græske Sprog som enten Plato eller
Aristoteles; men da Philologien fik sine Glarøine paa
og kom til at see efter, viste det sig at være en Feiltagelse,
og Orthodoxien maatte give efter. Saa skulde
da, trods menneskelige Ufuldkommenheder i det Sproglige,
hele den bibelske Historie fra Først til Sidst i Et
og Alt være ufeilbar. Men nu toge Undersøgelserne et
nyt Opsving. Lærde Forskere beviste, at mange af
Bøgerne vare uægte; at Texten var tvivlsom; at der
fandtes Uovereensstemmelser, ikke alene mellem de
bibelske Beretninger og profanhistoriske Vidnesbyrd,
men, hvad der endnu var værre, Afvigelser og øiensynlige
Modsigelser imellem de hellige Forfattere indbyrdes.
Tilsidst stege Indsigelserne til en saa overvældende
Flod, at den gamle Orthodoxie stod maalløs,
fortabt, og saae alle sine Dæmninger brudte.

Saa indtraadte der da for Kirkelærens Vedkommende
en mere liberal og fordomsfri Periode. Det
% --- p. 34 ---
\setcounter{footnote}{0}%
\opage{34}kunde jo ikke hjælpe at være halsstarrig og forsvare
det Uforsvarlige; man maatte bøie sig for Videnskabens
Beviser, man maatte være fornuftig. Men efterhaanden
som man selv begyndte at blive fornuftig, begyndte
man da ogsaa at indsee, at der i den hellige Skrift
var Meget, som stred mod Fornuften; Djævelen f. Ex.
var slet ikke fornuftig, Dæmonerne og Besættelserne og
alt det Overnaturlige, hvormed den skrevne Aabenbaring
var belemret, var ufornuftigt. Til Syvende og Sidst opdagede
man, at alle Mirakler --- naar man undtog
Skabelsens --- vare ufornuftige.

Kun det Naturlige var virkeligt; det Overnaturlige
var uvirkeligt. Man maatte da see at faae det Overnaturlige
bort; det kunde allerbedst skee ved en philologisk-kritisk
Fortolkning. Dog her stødte man paa en
Anstødssteen, der ikke var let at rydde af Veien; thi
det var hele Bygningens Hjørnesteen, det var Christus
selv. Man fandt det beklageligt, at en ellers saa fornuftig
Mand som Christus havde indladt sig paa at
gjøre Mirakler. Fritænkerne erklærede ham aabenlyst
for en Bedrager; men Kirkelærerne, der godt kunde
see, at han var et Dydsmønster, undskyldte ham paa
det Bedste og forsikkrede, at han havde gjort det i en
god Mening, blot for at læmpe sig efter Almuens Forestillinger.
Andre skøde Skylden fra Mesteren over paa
Disciplene; det var Disciplene, der havde misforstaaet
ham, og forvandlet hans naturlige Sygehelbredelser til
Mirakler. Og hvad blev der saa af den evangeliske
Historie? En tør Beenrad i forvrængende Omsvøb. Den
rene Jomfru blev til en ordinær Tømmermandskone, ja
til en Kvinde af tvetydigt Rygte. Den Opstandne var
en Skindød, der skrantende af sine Saar listede sig
% --- p. 35 ---
\setcounter{footnote}{0}%
\opage{35}omkring i Tusmørke og kun fik en glimrende Bortgang,
for saa vidt han engang i en stærk Aftenbelysning gik
bag et høit Bjerg og derved fik Disciplene til at troe,
at han var faret til Himmels. Tarveligt var det med
denne Aabenbaring; men det var videnskabeligt. Da
viste der sig med Et en ny Udvei, det var et nyt Syn;
Æsthetiken og Philosophien havde faaet Øje paa Phantasien,
at sige den menneskelige.

Hvorfra stamme de hedenske Guder? Det er ikke
historiske Personer, ikke tomme Naturallegorier; heller
ikke er det vilkaarlige Digtninger, som ved et Præstebedrag
have fundet Indgang hos det menige Folk og
grundlagt dets Overtro. De hedenske Guder ere Skabninger
af en ideeopfyldt, men ubevidst digtende Phantasie.
Disse Frembringelser i Sagn og Digt ere Myther;
Religion er Mythologie. Nu slog Befrielsens Time. Ere
de hedenske Religioner Mythologier, hvorfor skulde den
saakaldte aabenbarede Religion, i Jødedom og Christendom, ikke ogsaa være Mythologie? Mythen har netop
det Eiendommelige ved sig, at den med en aandig
Naturkrafts Uimodstaaelighed udspringer i Phantasien.
Den skabende Phantasie laaner sine Elementer fra
Naturbegivenheder og Historie; det er da heraf forklarligt,
at det Mythiske kan paa den allernaiveste
Maade være sammenblandet med det udvortes Virkelige,
det Overnaturlige med det Historiske. Nu begyndte
man da at finde Myther først i det Gamle og siden i
det Nye Testamente, Naturmyther, mythiske Sagn, historiske
Myther, philosophiske Myther o. s. v. Og hvad er
saa Resultatet? Det nyeste videnskabelige Resultat er,
at alt det Overnaturlige i det Nye saa vel som i det
Gamle Testamente er Mythologie. Aanden er rigtignok
% --- p. 36 ---
\setcounter{footnote}{0}%
\opage{36}borte, men man har jo dog Ideen. Paa Ideens Vegne
skal den officielle Religionslære for Tiden anvende Kritiken
for at udfinde, hvad og hvormeget der i Evangelierne
er Mythe og hvad der er virkelig Historie.

Jesu Liv falder i tre Afsnit: 1) fra hans Undfangelse
og Fødsel til hans Daab i Jordans Flod og hans
Fristelse i Ørken; 2) fra hans Optræden som Lærer til
hans Død og Begravelse; 3) det tredie og sidste Afsnit
omfatter hans Opstandelse og Himmelfart. Ifølge den
nugjældende videnskabelige Synsmaade indeholder kun det
andet Afsnit om hans offenlige Liv en virkelig Historie,
og fra denne Histories Virkelighed maa da hans Mirakler
udskilles, thi Miraklerne ere jo kun Myther. Det første
og sidste Afsnit er reen og skær Mythologie. Bebudelsen en Mythe; Fødselsunderet en Mythe; Opstandelsen med de paafølgende Aabenbarelser og Himmelfarten, en Mythe. Ingen, der vil staae paa Videnskabens
 Grund er i Stand til her at hævde det historisk Virkelige;
det er ikke een Ting eller to Ting, som her give
Vanskeligheder; det er, om man vil, tusinde, ja titusinde
Ting, thi det er alle nuværende Videnskabers samstemmende
Vidnesbyrd om Naturlovene og deres Uforanderlighed,
Forsvareren af det Historiske her faaer imod
sig. Det er lige saa umuligt, at faae nogensomhelst af
de mange Videnskaber overtydet om, at det ikke er
Myther, men virkelige Kjendsgjerninger, som det er
umuligt at overbevise Astronomien om, at Himlen er
en fast Hvælving, og at Jorden staaer stille i Verdens
Midtpunkt. Den Enkelte kan gjerne paastaae det; men
Videnskaben agter ikke hans Paastand. Troens Forsvarere
kunne beraabe sig paa, at Christendommen er
en historisk Kjendsgjerning; Videnskaben svarer, at
% --- p. 37 ---
\setcounter{footnote}{0}%
\opage{37}Muhamedanismen, Brahmaismen, Buddaismen ogsaa ere
historiske Kjendsgjerninger, uden at deraf kan udledes,
at deres Myther ere guddommelige Aabenbaringer.

Saa staaer da endnu den gamle, traditionelle, men
aandsforladte Religionslære paa Videnskabens Grund.
Men hvordan staaer den? Dens Stilling er ikke misundelsesværdig,
den seer med foldede Hænder paa de store
Aabenbaringsbilleder og gribes af Veemod. Bebudelsens
Engel, Hyrdernes Syn i den straalende Nat, hvor enfoldig ophøiet. Gid det var virkeligt! Men desværre,
det er ikke virkeligt, det kan ikke være virkeligt.
Videnskaben siger Nei! og Naturlovene true: saa faaer
det i Guds Navn at være Myther.

Dersom Sagen ikke var saa alvorlig, at Situationen
i Sandhed maa blive tragisk, saa vilde den hele Fortvivlelse
virkelig være komisk; thi den er selvforskyldt:
det Hele beroer paa en Indbildning. En Digterphantasie,
der vil gjøre sig bekjendt med den officielle Religionslæres
nyeste videnskabelige Bedrifter, vil finde et uudtømmeligt
Oplag af Emner til komisk Behandling. Snart
hedder det sig, at Johannesevangeliet er det sande dybe
Hjertes Evangelium og at de tre andre maae rette sig
efter Johannes. Saa kommer der med næste Post Efterretning
om, at Johannesevangeliet er et fiint udarbeidet, bevidst
Konstværk. Altsaa ikke historisk! Hvem af Evangelisterne
skal nu spille Hovedrollen? Lukas kan det nu
ikke være: altsaa enten Mathæus eller Markus. Det er
Mathæus; nei, det er Markus! Saa tænder man Kritikens
Løgte og leder efter et Urevangelium.

Lige saa stor er Raadvildheden med det Historiske
i Miraklerne. Man vil skjelne kritisk mellem
Virkelighed og Sagn, og kun antage det for Virkeligt,
% --- p. 38 ---
\setcounter{footnote}{0}%
\opage{38}som har alle Evangelisters samstemmige Vidnesbyrd.
Vandringen paa Havet behøver da ikke at være virkelig,
--- at sige paa Kathedret, thi paa Prædikestolen er den
altid virkelig; men Bespiisningsmiraklet maa nødvendig
være virkeligt, thi det har jo alle Vidnesbyrd. Dog nei,
det kan umulig være virkeligt: Brødene ere for faa og
 Fiskene for smaa. „At see Brødet svulme og formere
 sig i Christi Haand, er et uhyggeligt Syn: hvad vil
 Videnskaben sige?“ --- Var det saa ikke bedst at ombytte
det videnskabelige Grundlag med et religiøst, og vælge
 sit Standpunkt udenfor Videnskaben? „Nei! Nei! det
gaaer paa ingen Maade an; den Dag da Religionslæren
ophører at være en Videnskab, en objectiv, paa objective
Kjendsgjerninger fast begrundet Videnskab, paa
den Dag maa Religionen gaae tilgrunde, og det sørgelige Budskab vil staae i Avisen. Og det behøves heller
ikke; Faren er ikke saa stor; det er kun en Overgang.
Videnskaben har endnu ikke sagt sit sidste Ord. Man
kan være vis paa, at naar Videnskaben en Gang faaer
sagt sit sidste Ord, saa skulle vi nok faae Miraklernes
historiske Virkelighed objectivt videnskabelig begrundet.“
Men naar faaer Videnskaben sagt sit sidste Ord? Det
bliver en langvarig Proces, længere endnu end den fra
Bleakhouse bekjendte: Jarndyce and Jarndyce. Den
menneskelige Forstand er forarget, og Phantasien forvildet;
den af Forargelse forvildede Phantasie staaer
smilende mellem Actstykkerne og siger med den gamle
vanvittige Jomfru: Pray, accept my blessing!

Religionslæren behøver ikke at gaae ud af Videnskabernes
Kreds; men den behøver i Sandhed at stille
sig paa et ikke-videnskabeligt d. e. paa Troens og Aaben%
% --- p. 39 ---
\setcounter{footnote}{0}%
\opage{39}baringens eget overvidenskabelige Standpunkt. Lad saa
kun Videnskabens Lys falde opklarende over det Hele.

Stille vi os paa Religionens eget oprindelige Grundlag,
vil det strax blive indlysende, hvorledes det gaaer
til, at Mythe og Aabenbaring kunne have en saa paafaldende
Lighed med hinanden, at den videnskabelige
Kritik, der har nedsat det Overnaturlige til Skin og
Phantasie, ikke, trods al sin Skarpsindighed, er i Stand
til at skille dem ad. De have begge en fælles Form,
Synets, Phantasiens Form. Den almægtige Gud er ogsaa
Synets og Synernes Gud. Men der er en Forudsætning,
hvorpaa Troen staaer fast, og denne Forudsætning
kan ingen Videnskab rokke, det er Kjærlighedens Almagt.
Videnskaben kan hverken benegte eller bekræfte
Guds Kjærlighed; thi den veed slet Intet om den. I
Kjærlighedens Almagt ligger Aabenbaringens Mulighed;
denne Mulighed skal Videnskaben vel lade staae, og
saa er det nok. I Aabenbaringens Mulighed har Troen
et Skjold mod Kritikens Pile; men i Virkeliggjørelsen
af denne Mulighed er ogsaa Phantasiens Magt aabenbaret.


Kjærligheden er en Tanke og mere end en Tanke;
den er en Villie og mere end en Villie; den har et
Hjerte, og Hjertet begjærer et Syn, og i Synet en Skikkelse,
og Syn og Skikkelse sees i Phantasie. Der er i
Mythen Noget af en guddommelig Phantasie; men det
Guddommelige er fordunklet. Hildet i Sandselighed, besnæret
af selvisk Attraa, har den til sig selv overladte
Menneskephantasie, kun seet Glimt af Sandheden, ikke
Sandheden selv. Lyset har skinnet i Mørket, men
Mørket har ikke begrebet det. Derfra det Forblommede,
det Uklare og Tvetydige; derfra det phantastisk
% --- p. 40 ---
\setcounter{footnote}{0}%
\opage{40}Illusoriske, hvormed Mythen altid er behæftet. I Aabenbaringen derimod har den guddommelige Phantasie Magten over den menneskelige, klarer den i sit Lys og
fylder den med sin Sandhed; derfor ere Aabenbaringsphænomenerne
saa rene, utvetydige, gjennemsigtige for
Aanden. Aand og Natur ere inderlig forenede; ogsaa
det Overnaturlige har sin Natur. Naturbegrebet har
derfor tvende Sider, en videnskabelig og en religiøs.
Det videnskabelige Naturbegreb dannes og maa danne
sig for os under Indtryk og Iagttagelse af naturlige
Kjendsgjerninger; men gjennem det Synlige skimtes det
Usynlige, gjennem det Bekjendte det Ubekjendte, gjennem
Lovene Lovgiveren, gjennem Nødvendigheden Friheden.
Aanden i sin Uforanderlighed med sine Evighedstanker
udfolder Naturindholdet i Tingenes Sammenhæng,
i Phantasienødvendighedens Form. Længere kommer
Videnskaben ikke.

Hvor Videnskaben slipper, begynder Religionen, og
Religionen er det Oprindelige; den gaaer i Tiden forud
for al Videnskab; den gaaer igjennem Aabenbaring og
Myther. Her gjælder det da om at gjøre Skilsmisse
mellem Indbildning og Sandhed, at finde en Prøvesteen,
der kan adskille Aabenbaring fra Mythe. Denne Prøvesteen
have vi i Troen paa Kjærlighedens Almagt. At
Naturens Skaber ogsaa er Aandernes Fader, at Nødvendighedens
og Lovenes Gud ogsaa er Frihedens og
Kjærlighedens Gud, at det Overnaturlige i fri Inderlighed
er det usynlige Grundlag for alt det udvortes Naturlige,
det kan kun erkjendes i Tro. Denne Tro, een
Gang plantet i Slægtlivets Grund og udtalt i en bestemt,
utvetydig Troesbekjendelse, faaer Videnskaben
aldrig udryddet af Menneskenes Hjerte. Ramsen om
% --- p. 41 ---
\setcounter{footnote}{0}%
\opage{41}de uforanderlige Love bliver  intetsigende, naar det erkjendes,
at der ikke blot er en synlig Sammenhæng
mellem de synlige Ting, men ogsaa en indre, om end
usynlig Sammenhæng mellem det Naturlige og det
Overnaturlige, at, med andre Ord, Verdener af lavere
og høiere Orden ere forenede i Almagtens Væsen, i Frihedens
og Kjærlighedens guddommelige Phantasie.

Forskjellen mellem Aabenbaring og Mythe kan
sees i Troens Lys paa Tanken, Villien og Phantasien.
Paa Tanken! Der er Tanker i Mytherne, ofte dybsindige
Tanker, men Synerne gaae op og ned, de krydsende
Tanker blandes, og Forestillingerne forvirres. I
Aabenbaringen er Tanken klar, Lovens Tanke, Livets
Tanke, Tanken om Synd og Død og Tanken om Opreisning
og Frelse. Paa Villien! Gjennem Mytherne gaaer
der en Længsel, en Higen og Trang til at see en nærværende
Gud. Men Sandseligt og Aandeligt, Naturdrift
og selviske Lidenskaber blande sig med den dybere
Attraa, og den hellige Villie vedkjender sig det ikke.
Saa blive Mythens Undere Indbildninger, phantastiske
Tegn: her skee ingen Mirakler. I Aabenbaringens Ord
er Villien med, den almægtige Villie: Tanke og Villie
svare til hinanden som Ord og Tegn. Aandens Gud,
Lovens og Kjærlighedens Gud er med i sit Ord og
taler i Tegnet Almagtens Sprog: Hører Ordet, see her
mit Tegn! det er Miraklet. Men saa udvortes virkeligt
end Miraklet kan være, det maa dog sees med
Troens Øje: saa sætter den guddommelige Phantasie den
menneskelige i Bevægelse. Hvilken Forskjel imellem den
ofte glimrende, farverige, men trods det Storslaaede og
Plastiske dog forvirrede Mythephantasie og den klare
af Salighedslyset bestraalede, og just derfor saa rene,
% --- p. 42 ---
\setcounter{footnote}{0}%
\opage{42}kydske Aabenbaringsphantasie. Kaste vi et Blik paa
Bebudelsens Evangelium og lade den guddommelige
Tanke og Villie og Phantasie være forenede i Kjærlighedens
Almagt, saa kunne vi forstaae Aabenbaringen.
Her er en Tanke, der skal indføres i Menneskelivet, en
Kjærlighedstanke, og Tanken er denne: da alle selv
de ædleste Naturer i Slægten bære Spiren i sig til
Synd og Død fra Fødselen af, saa skal der nu fødes
Een, som af Naturen selv er usyndig og reen, et Aandens
Menneske, et Evighedens Barn, en hellig Spire,
som voxer op uden Synd og frelser fra Synd og Død.
Denne Tankes Eenhed med den almægtige Villie er
Underets Aarsag. Underet selv er en Hemmelighed;
men Hemmeligheden skal aabenbares, den skal blive
vitterlig for Alle; saa kommer den Alt aabenbarende
Phantasie og gjør det Usynlige synligt. Englenes Hilsen,
Marias Ord, den høitidelige Erklæring! i dette
Livsbillede er der intet mythisk, intet tvetydigt; det
siger netop, hvad det skal sige, hverken mere eller
mindre. Det er ikke menneskelig Indbildning; det er
i Kjærlighedens Almagt guddommelig Phantasie; hver
Linie og hver Farvetone, hver enkelt Træk og hver
Lysning derover, Alt er af Aand; det saaledes Aabenbarede
er virkeligt i allerhøieste Forstand, men det er
Aandsvirkelighed, Virkelighed for Aand. --- Som Begyndelsen
er, saa er ogsaa Fortsættelsen og den sidste Slutning.
Men Lysningerne vexle, det Naturlige og det Overnaturlige,
det virkelige Udvortes og det høiere Indvortes
gribe vidunderligt ind i hinanden; kun Aabenbaringsphantasie
i Eenhed med Tro formaaer at klare os en saadan
Vexelvirkning: vi skulle nærmere see, hvorledes.

% --- p. 43 ---
\setcounter{footnote}{0}%
\lecture{Fjerde Forelæsning.}\opage{43}

Hvor let en Strid om guddommelige Ting kan udarte
til en tom Ordstrid, naar der ikke tilbørlig sørges for,
at Ord og Sag gjensidig belyse hinanden, er nok som
bekjendt. Strides der f. Ex. om Tro og Viden, kan
man overmaade let give det Skin af, at al Viden er
grundet paa Tro. Maleren, der skal bruge Farverne,
maa jo troe, troe paa sit Syn; Mathematikeren med sit
a = a maa troe, troe paa sin Forstand; Astronomen
maa troe paa sine egne og Andres Iagttagelser o. s.
frmdl. Men man kan da ogsaa fra den anden Side
ligesaa let give det Skin af, at al Tro er grundet paa
Viden. Vi maae jo dog vide, hvad Tro er, vide, hvad
det er, hvorpaa vi troe, og vide, hvorfor vi troe derpaa;
og naar vi vide Alt dette tilgavns, er det Hele jo bygget
paa Viden. Dette Spilfægteri hører op, naar vi ærlig
gaae ind paa Sagen. Der er Noget, der ikke skal troes,
men vides, og der er Noget, der ikke skal vides, men
troes. Siger man, at Solen skinner, saa svarer den
Seende: Jeg veed det; kun en Blind kunde sige: jeg
troer det. Siger man at to og to er fire, saa svarer
den Forstandige: jeg veed det; kun en Idiot kunde sige:
jeg troer det. Er Talen om en Pengesum, og der
% --- p. 44 ---
\setcounter{footnote}{0}%
\opage{44}spørges; er Regningen rigtig? da hjælper det ikke at
svare: jeg troer det, thi det skal ikke troes, det skal
vides. I slige Forbindelser betyder Tro kun en foreløbig
Antagelse, en Halvvished forenet med Tvivl. Her
søges Viden, først med Viden kommer Visheden, og med
Viden og Visheden forsvinder Troen. Den, som i slige
Tilfælde har naaet til fuld Vished, siger: jeg troer det
ikke, thi jeg veed det. Anderledes, naar Talen er om
Sjælens Udødelighed, om Guds Kjærlighed, om den overnaturlige
Verden. Et Menneske kan, hvad alt dette angaaer,
være saa vis i sin Sag, at han uden Omsvøb tør
sige: jeg veed det; jeg veed det vist, at Gud er Kjærlighed;
jeg veed det vist, at der er et evigt Liv, en
Opstandelse o. s. v. Viden betyder her en urokkelig
Vished; men denne urokkelige Vished gjør ikke Troen
overflødig, tvertimod! den er Et med Troen selv: Viden
er her et Udtryk for den fulde Forvisning i Tro. Ogsaa
fra dette Synspunkt kan Tro sættes imod Tvivl,
men Tvivlen har da ikke sit Udspring fra ufuldkommen
Viden, men fra Vantro. Det, som i alleregenligste Forstand
maa troes, er netop det, som ifølge sin Natur
ikke kan vides; --- kun maa det ikke oversees, men paa
det Eftertrykkeligste indskærpes, at Troen kan i sin
Erkjendelsessphære give en lige saa tilforladelig Vished,
en lige saa fuldkommen Forvisning som Viden i sin.

I Troens Sphære er Modsætningen mellem Tro og
Viden forsvunden, thi her spørges, i videnskabelig Forstand,
slet ikke om Viden; her indtræde Modsætninger
af en ganske anden Art: Modsætning mellem Tro og
Overtro, mellem Vantro og Overtro.

For at komme til Klarhed over, hvorledes det forholder
sig med Modsætningerne paa Troens Omraade, er
% --- p. 45 ---
\setcounter{footnote}{0}%
\opage{45}der et Forhold, vi igjen maae tage i Øjesyn, Grundforholdet
mellem Phantasie og Virkelighed.

Al Virkelighed, den være udvortes eller indvortes,
sandselig eller oversandselig, falder --- ikke udenfor ---
men indenfor den almægtige Phantasie. I denne allerhøieste
Forstand kunne vi da sige, at al Virkelighed er
Aandsvirkelighed, den har jo i Phantasien sit Udspring
fra Aandens skabende Almagt.

Men fra Menneskets Standpunkt seer Sagen anderledes
ud. Her er en Adskillelse, en Kløft imellem en
udvortes og en indvortes Virkelighed, imellem Tingene
udenfor os og vor Bevidsthed om Tingene. Tingene
maae sandses; vel er der i Sandsningen en skjult
Phantasie, for saa vidt der er en Anskuen; men denne
Phantasie er bunden i Sandsningen, og Sandsningen,
bunden og betinget af Indtryk fra Tingene, er umiddelbart
afhængig af det udvortes Givne. --- Dette giver
os en udvortes Virkelighed, Naturvirkelighed, Virkelighed
for Sands og Forstand. I Modsætning hertil have
vi en indvortes Verden, en Verden af Billeder, Forestillinger
og Tanker, en oversandselig Verden, en Verden,
hvori ogsaa det Overnaturlige kan fremstille sig i Phantasiens
Medium. Denne Verden har Virkelighed, for
saa vidt den har Sandhed, og den har Sandhed, for
saa vidt den er af guddommelig Oprindelse. Hvor
Phænomener af denne oversandselige, overnaturlige Verden
da her træde os imøde med den Aabenbaringens Klarhed
og Kraft, der underlægger sig det Udvortes, der
tale vi i særlig og religiøs Forstand om Aandsvirkelighed,
Virkelighed for Aand.

Den udvortes Naturvirkelighed imponerer det sandselige
Menneske og staaer for Bevidstheden med Præg af
% --- p. 46 ---
\setcounter{footnote}{0}%
\opage{46}en Selvstændighed, der sætter Phantasien Skranker og
forvandler dens Billeder til Indbildninger; men det
aandelige Menneske erkjender, at alt det Udvortes har
sin Oprindelse fra det Indvortes, at Naturvirkeligheden
med dens Kjendsgjerninger og Love er og maa være
Aandsvirkeligheden absolut underordnet.

Fra dette Stade kunne vi nu med særligt Hensyn
paa Phantasien allerede see, hvorledes det gaaer til, at
Vantro og Overtro stille sig mod Tro. Naar Sandser
og Sandselighed i den Grad har fanget og fængslet et
Menneskes Phantasie, at det Naturlige alene har Gyldighed
for ham, medens det Overnaturlige er blevet
forvandlet til Skin og Indbildning, saa have vi Vantroen.
Vantroen kjender ingen anden Phantasie end den blot
menneskelige, den æsthetisk illusoriske. Har omvendt
det Overnaturlige i Phantasien antaget Charakteer af
et udvortes Sandseligt, saa at det Aandsvirkelige sættes
paa lige Trin med det i Naturen Materielle og Virkelige,
saa have vi Overtroen.

Kun Troen formaaer at opfatte begge Virkelighedsformer
i deres indbyrdes Forskjel og Eenhed; thi kun
Troen formaaer at erkjende en Vexelvirkning mellem
den guddommelige og den menneskelige Phantasie. Der
er i det Gamle Testamente storstilede Exempler, som
lære os, hvad dette vil sige.

Moses taler som Lovgiver aldrig i sit eget, men
altid i Jehovas Navn. Forkyndelsen af de ti Bud begynder
med den bekjendte Formel (2 Moseb. 20):
„Gud talede alle disse Ord og sagde.“ For den religiøse
Sands kan dette i sin Almindelighed ikke have noget
videre Stødende. De ti Bud ere af almeensædeligt
Indhold; de udtale Tanker og Grundsætninger, som
% --- p. 47 ---
\setcounter{footnote}{0}%
\opage{47}ethvert borgerligt Samfund overalt maa vedkjende sig:
man kunde næsten kalde dem Fornuftbud. Den Smule
Phantasie, der hører til for at sætte sig i Mose Sted,
erkjende hans Tro paa Jehova og indsee Nødvendigheden
af, at han maatte tale i Jehovas Navn, naar han
skulde tale med Myndighed, indfinder sig som af sig
selv. Det Overnaturlige og det Naturlige, det Overfornuftige
og det Fornuftige synes her uvilkaarlig at
smelte sammen til fuld Harmonie. Men Eenhedsbaandet
brister snart, naar vi gaae lidt videre i Texten. Det
hedder (2 Moseb. 25): „Og Jehova talede til Moses,
idet han sagde: Tal til Israels Børn, at de skulle tage
mig en Offergave af hver Mand\dots{} Og dette er Offergaven,
som I skulle tage af dem: Guld og Sølv og
Kobber. Mørkeblaat og Purpurrødt og Carmoisinrødt
og fiint Linned og Gedehaar og Sælhundeskind og rødfarvet
Vædderskind og Acacietræ. Olie til Lysning,
Urter til Salveolie og til vellugtende Røgelse\dots{}“ Er
dette ordret fra Jehovas Mund, saa gaaer Jehova med
sine Forskrifter for Tabernaklets Indretning unegtelig
vidt i det Smaalige; han indlader sig paa de allerubetydeligste
Enkeltheder. Ingen menneskelig Bygmester
kan gjøre det nøiagtigere. Det Ophøiede synes
her at blive borte i lutter Endelighed. Hvad skulle vi
nu fra Troens Standpunkt tænke og dømme? Tage vi
Forskrifterne som ordret dicterede af Jehova selv, saa
gaaer det Hele over i udvortes Virkelighed og bliver
en udvortes Historie. For Tilegnelsen af en saadan
udvortes Historie har den religiøse Tro ingen Prøvesteen,
intet indre, ubedrageligt Sandhedsmærke. Forskjellige
Religionsbøger, Zendbøgerne hos Perserne,
Vedabøgerne hos Hinduerne, Koranen o. fl. kunne jo
% --- p. 48 ---
\setcounter{footnote}{0}%
\opage{48}indeholde lignende Forskrifter. Credat Judæus Apella,
non ego. Men naar det indre, ubedragelige Sandhedsmærke
forsvinder, forvandler Troen sig til Overtro.
Overtroens Phantasie er kjendelig paa, at det Aandelige
nedsænkes i det Sandselige --- Gedehaar og Sælhundeskind
og rødfarvet Vædderskind kan kun beskjæftige
Phantasien paa en aldeles sandselig Maade --- saa
at Sandhedskriteriet gaaer tabt. Gaaer man nu til den
modsatte Yderlighed og paastaaer, at de særlige Forskrifter
for Tabernaklet ikke hidrøre fra Jehova, men
alene fra Moses, saa kan man ikke uden Vilkaarlighed
antage, at de ti Bud ere Jehovas Bud, thi med Forskriften
for Sabbaten og Forbudet mod at begjære, hvad
Næstens er, gaaes der jo ogsaa i det Enkelte --- Næstens
Hustru, hans Tjener, hans Tjenestepige, hans Oxe, hans
Asen kunne jo ogsaa beskjæftige Phantasien paa sandselig
Maade. Men naar det med Jehovas Ord svinder
ind til en Indbildning, saa have vi Vantroen.

Hvor er nu Troen med sit Sandhedsmærke? I
Vexelvirkningen mellem den guddommelige og den
menneskelige Phantasie; thi heri indesluttes den religiøs-psychologiske
Vexelvirkning mellem den guddommelige
og den menneskelige Tanke, den guddommelige
og den menneskelige Villie, det guddommelige og
dog menneskelige Syn, det guddommelige og dog menneskelige
Sprog i guddommelige og dog menneskelige
Billeder. Det lovgivende Ord er Ord af Moses og dog
i Sandhed Jehovas Ord; Billedet af Tabernaklet i alle
dets Enkeltheder er en Frembringelse af Mose Phantasie,
dannet, formet og udarbeidet i Prophetens Indre,
og dog et guddommeligt Grundsyn, Planen i Et og Alt
Jehovas Plan. Saa taler Israels Gud, og Talen er
% --- p. 49 ---
\setcounter{footnote}{0}%
\opage{49}virkelig, aandsvirkelig. Jehova taler til sin Prophet
som Aand til Aand, han taler til Moses, i Mose Indre
--- igjennem Moses, et guddommeligt og dog menneskeligt
Ord; saaledes blive Budene Mosebud, og dog Jehovas
Bud.

Hvad vide vi da om denne Vexelvirkning? Vi vide,
i videnskabelig Forstand, Intet, slet Intet; her er ingen
Gjenstand for Viden, men i Kraft af den religiøse Phantasie,
er her en Gjenstand for Tro. Sandhedsmærket
er ubedrageligt, naar Aabenbaringstanken er i Eenhed
med Troestanken; den aabenbarede Villie i Væsen
og Grund er Troens Villie: Sandhedsprøven er i det
troende Sind, og Forvisningen i det Indre.

Lad os gaae tilbage til Begyndelsen og betragte
Prophetens Kaldelse. I Midians Ørk ved Bjerget Horeb,
fjernt fra Pharaos Hof, er Moses Hyrde. Derude
mellem Bjergene er der en Eensomhedens Uendelighed:
Minderne vaagne, Tankerne tale, men Stemmerne tie,
Taushed boer i de øde Fjelde, derude er der et Møde
mellem Phantasie og Virkelighed. Og det Virkelige er
et Udvortes, en Tornebusk, og see! der er Ild i Tornebusken,
og Tornebusken brænder og bliver dog ikke
fortæret. Saa høres en kaldende Røst; det taler i
Moses til Moses: „Jeg er din Faders Gud, Abrahams
Gud, Isaks Gud og Jakobs Gud. Jeg har visselig seet
mit Folks Elendighed. --- Gaa, jeg vil sende dig til
Pharao, at udfrie mit Folk af Ægypten.“ Og Moses
skjælver ved Tanken om det store Hverv, der er uendeligt
over hans Kræfter. „Men“ --- lyder Stemmen ---
„hvad er det du har i din Haand?“ En Stav. „Kast den
paa Jorden“ --- og han kastede den til Jorden, og den
blev til en Slange, og Moses flyede tilbage for Slangen.
% --- p. 50 ---
\setcounter{footnote}{0}%
\opage{50}--- „Indbildning, Skin og Phantasie, et oprørende Brud
paa Naturlovene!“ Det er Stemmer fra Videns Kathedre;
men disse Stemmer og disse Raab, hvad enten de
komme fra de Lærde selv eller fra deres Skoledrenge,
naae ikke ud til hine Egne, forstyrre ikke Eensomhedens
Fred i Midians Ørk, ikke den høitidelige Taushed,
der holder Vagt om Bjerget Horeb. Det Eneste,
hvorom her spørges, er Aabenbaringstanken i dens
Eenhed med Villiens og Phantasiens Almagt. Den
Slange, som fremkommer af Mose Stav, har ingen Plads
i Zoologien; at stirre paa den som paa et udvortes
Sandseligt, er skadeligt for Øinene, man bliver blind
enten af Overtro eller af Vantro. Men, sige de Vidende,
her menes jo dog en virkelig Slange; Moses seer den:
han gyser og farer tilbage. Vi svare: ja, vel er den
virkelig, men Virkeligheden er af en egen Art, og maa
sees med Troens Øine. Det Overvirkelige er som et
Lyn i Mørket: i samme Nu som det sees, har Øiet
lukket sig; intet Skin er saa blændende som Skinnet
af det Overnaturlige, intet Lyn er saa stærkt som Almagtens
Udslag.

Det er just det Eiendommelige ved et saadant
Aabenbaringssyn, at det har Skikkelse af et Udvortes
og ligner det Naturlige, og dog, naar Sands og Forstand
ville til at undersøge det, saa er der enten et
blot Naturligt, eller det Hele er forsvundet. Kun Troens
Øie udholder at see paa Tegnet, thi dette Sjælens Øie
er netop bygget og organiseret saaledes, at det Sandselige
kan opløse sig i et Oversandseligt, og det Udvortes
omsættes i et Indvortes. Og hvor er saa Sandhedens
Mærke? I Tankens Eenhed med Villiens Almagt,
i den guddommelige Phantasie, der gjør det Usynlige
% --- p. 51 ---
\setcounter{footnote}{0}%
\opage{51}synligt. Jehova vil frelse sit Folk fra Ægyptens Plager,
fra Trældommens Aag, han vil frelse det; her er hans
Ord. Han kan frelse det, her er hans Tegn! Ja, her
er han selv: med denne Forvisning, en Troesforvisning,
der holder Stand mod al Videns Tvivl, tager da Moses
sin Stav op og gaaer til Pharao.

At det udvortes Naturlige kan i Aandsvirkeligheden
gribes og magtes af det Overnaturlige, er umiskjendeligt;
men at det i den religiøse Phantasie ogsaa kan
virke paa fri indvortes Maade, er klart af de prophetiske
Syner. Vi vælge et Sted hos Esaias (6. Kap.).
Propheten beskriver sin Kaldelse. „I det Aar, Kong
Usias døde, da saae jeg Herren siddende paa en høi
og ophøiet Throne, og hans Klædebons Flige fyldte
Templet. Serapher stode over ham, hver med sex Vinger;
med de to tildækkede hver sit Aasyn, med de to
tilhyllede han sine Fødder og med de to fløi han. Og
den ene raabte til den anden: hellig, hellig, hellig er
Jehova Zebaoth, hele Jorden er fuld af hans Herlighed!
Og Dørposterne skjalv ved den Raabendes Røst og
Huset fyldtes af Røg. Og jeg sagde: vee mig! jeg
gaaer tilgrunde; thi jeg er en Mand med urene Læber
og jeg boer midt iblandt et Folk, som har urene Læber,
og dog have mine Øine seet Kongen, Jehova, Zebaoth.
Da fløi en af Serapherne til mig, og i sin Haand havde
han en glødende Steen, hvilken han tog med en Tang
fra Alteret. Og han rørte ved min Mund og sagde:
see, dette berører dine Læber; og din Misgjerning er
vegen fra dig, og din Synd skal forsones. Og jeg
hørte Herrens Røst, som sagde: hvem skal jeg sende,
og hvem vil gaae for os? Da sagde jeg: her er jeg,
send mig!“ --- Dersom dette ikke er Aabenbarings%
% --- p. 52 ---
\setcounter{footnote}{0}%
\opage{52}phantasie, og dersom der i denne Aabenbaringsphantasie
ikke er en sand og levende Vexelvirkning mellem det
Guddommelige og det Menneskelige, hvor er der saa et
Sandhedens Lys? Er et Syn med dette Høihedspræg i
denne Symbolik med denne Reenhed ikke Virkelighed,
Aandsvirkelighed, Virkelighed for Aand: hvad er saa
Virkelighed og hvad er Aand?

Hvorledes den religiøse Phantasie, uden at tabe sit
Aabenbaringspræg eller opgive sin Troescharakteer,
gaaer over til Poesi, kan sees af mangfoldige Exempler
i det Gamle Testamente, navnlig hos David. Det Eiendommelige
ved denne Poesi er, at enten den vender
sig ind efter, klagende i Synd og Anger, eller den
vender sig ud efter og optager Indtryk fra den omgivende
Natur, beundrende dens Storhed, jublende over
dens Herlighed, overalt er Phantasien fyldt af Jehovas
Nærværelse og seer i Alt en Aabenbaring af Skaberens
Almagt. David skildrer, hvorledes Jehova i Farens
Stund kom ham til Hjælp mod hans Fjender, og aabenbarede
sig i et Uveir. „I min Angst paakaldte jeg
Jehova, og han hørte fra sit Tempel min Røst. Da
blev Jorden rystet saa den bævede, og Bjergenes Grundlag
skjalv.“ --- Det er i et Tordenveir at Jehova kommer.
„Der slog Røg op i hans Næse, fortærende Ild
foer ud af hans Mund, glødende Kul gnistrede fra ham.
Han bøiede Himlene og steg ned; han red paa en
Cherub og fløi, han svævede paa Veirets Vinger.“
Denne Fremstilling tyder ikke paa noget særeget
Gjennembrud af det Overnaturlige; det er den menneskelige
Side af Troesphantasien, som her er i fri Bevægelse;
Fremstillingen er saa klar, at vi paa alle Punkter
kunne see, hvorledes det udvortes Naturlige forvandler
% --- p. 53 ---
\setcounter{footnote}{0}%
\opage{53}sig i Speilinger af Digterkongens religiøse Phantasie.
Men denne Phantasie vilde ikke i Sandhed være religiøs,
dersom den ikke indenfra var grebet af Jehova selv og
beaandet af Aabenbaringens Indhold.

Hvor dybtgaaende Vexelvirkningen mellem det Guddommelige
og det Menneskelige er i det Gamle Testamentes
Aabenbaringsphantasie, viser sig endnu fra en
anden Side. Der er en væsenlig Forskjel mellem de
mythiske Guder, de græske f. Ex. eller de nordiske, og
Aabenbaringens eneste sande Gud. De mythiske Guder
ere sjælelig-legemlige Individualiteter, Menneskevæsener,
kun af høiere Orden. Der er menneskelige Tilbøieligheder
og Lidenskaber i disse Guder. De skifte Sind,
de blive vrede og overile sig; det gaaer saa vidt hos
Homer, at Guderne alene ved at see paa Kampen for
Troja og tage Parti blive saa indbyrdes splidagtige og
paa det Sidste saa iltre, at de bogstavelig slaaes saa
det ryger om Ørene.

Jehova har ingen saadan begrændset Individualitet;
han er uendelig ophøiet, uforanderlig; Jehova er Aand,
hans Væsen det evige, alvidende, selvmægtige Jeg. Da
Moses spørger, hvad Navn han skal give ham, naar han
kommer til Ægypten og taler til Israels Børn, faaer
han til Svar: Jeg er den jeg er, mit Navn er: Jeg
skal være. Og dog stiger den Ophøiede fra sin Evighedsthrone
ned til Jorden, den Uendelige indlader sig med
det Endelige, den Usynlige gjør sig synlig, ja hvad mere
er, han faaer et menneskeligt Sind, hans Stemninger
vexle, hans Barmhjertighed kan omskiftes til flammende
Vrede, han kan angre, fortryde; til hvilken Yderlighed
han kan gaae i sin Forbitrelse, naar hans Hæftighed
blusser op, derpaa have vi et ret slaaende Exempel i
% --- p. 54 ---
\setcounter{footnote}{0}%
\opage{54}anden Mosebog (Kap. 32). Da Aaron har støbt Guldkalven,
siger Jehova til Moses: Jeg har seet dette Folk,
og see, det er et haardnakket Folk. Og nu, lad mig
være, at min Vrede kan optændes mod dem og jeg kan
fortære dem, og jeg vil gjøre dig til et stort Folk. Da
bønfaldt Moses for Jehovas, sin Guds Aasyn og sagde:
hvorfor er din Vrede blusset op imod dette Folk, hvilket
du har udført af Landet Ægypten med en stor
Kraft og en stærk Arm? Hvorfor skulle Ægypterne
sige saaledes: til deres Ulykke har han ført dem ud,
forat slaae dem ihjel paa Bjergene og ødelægge dem
fra Jordens Overflade, vend om fra din opblussende
Vrede og fortryd dog det Onde imod dit Folk. Mindes
Abraham og Isak og Israel, dine Tjenere, hvilke du
har tilsvoret ved dig selv, til hvem du har sagt: jeg vil
mangfoldiggjøre eders Sæd som Stjernerne paa Himmelen,
og alt dette Land, som jeg har sagt, vil jeg give
eders Sæd, og de skulle arve det evindelig. „Da fortrød
Jehova det Onde, som han havde sagt, at han
vilde gjøre imod sit Folk.“ Det er en høist charakteristisk
Situation. Vantroen hoverer: her kan man see,
at Aabenbaringstro er Overtro og Overtroen dum. Propheten
overgaaer sin Gud: Moses er roligere, sindigere
end Jehova selv. Det er med ophøiet Overlegenhed, at
den høitbetroede Tjener bringer sin kongelige Herres
oprørte Sind i Ligevægt. Og hvorved skeer det? Ved
at vække den Forbittredes Æresfølelse; thi Ægypterne
vilde jo lee ad ham, naar han opgav sin Plan og slog
dem ihjel, han skulde befrie. Men især ved at minde
ham om hans Ed, at han ikke skal staae som en Meeneder
o. s. v.

% --- p. 55 ---
\setcounter{footnote}{0}%
\opage{55}Og hvad er saa det, som her forarger? Aandens
hellige Vrede.

Men med denne hellige Vrede har det i den religiøse
Phantasie en egen Sammenhæng. Aabenbaringens
Gud har et Hjerte, et Sind. Han kan elske og han
kan hade, han kan vredes og han kan formildes. Idet
hans Tanker, hans Villie og Synerne af hans Phantasie
gaae over i Menneskets Tanker og sætter Menneskets
Sind og Lidenskaber i Røre, da see vi den hellige
Vrede bruse op som et Menneskes Vrede og den guddommelige
Pathos gløde med en Ild af menneskelig
Lidenskab. Vi maae aldrig glemme, at det, der taler
til Moses, ogsaa taler i Moses. Sandhedsmærket sees,
naar der skjelnes mellem Tro og Overtro. For Overtroen
er hele Historien udvortes; hvad Jehova siger og
hvad Moses siger staaer udvortes adskilt, som naar det
ene Menneske taler til det andet; saa kommer Vantroen
og paaviser, at det Hele er en Mythe. Men Troen har
en anden Forklaring; den giver Jehova Ret, ubetinget
Ret, at det Folk, der efter en saadan Befrielse kan
dandse om Guldkalven, har fortjent Døden; den takker
Moses som Forbeder, men den seer tydelig at Henviisningen
til Ægypternes Hoveren og til den svorne Ed
ikke sigter paa en Tilstandsforandring i den Eviges
uforanderlige Væsen, men paa en Bevægelse i den
Bedendes Indre. Det er Prophetens Tro paa den frelsende
Jehovas uforanderlige Villie, der løfter ham op
over hans egne menneskelige Lidenskaber. Troen forstaaer,
hvad det vil sige, at den evig høie Aandens Gud
faaer et menneskeligt Præg, ja et folkeligt Præg.

Israels Gud kar et jødisk Physiognomie; han sees % PRINTED AS IS: „kar“ for har
ikke altid i det rene, klare Himmellys, men i Blandings%
% --- p. 56 ---
\setcounter{footnote}{0}%
\opage{56}skjær af jødisk Folkesind, et Luehav af jødiske Lidenskaber.
Det er ikke for Intet, at han kalder sig Abrahams
Gud, Isaks Gud og Jakobs Gud. Med dette
Physiognomie er han den Ophøiede, hans Aasyn mildt,
og dog, naar Vreden kommer, forfærdeligt.

Aabenbaringsphantasien i det Gamle Testamente
er jødisk: rig paa Forjættelser er denne Jehovaphantasie,
mægtig overvældende i sine Trusler. „Jeg
vil hjemsøge Jerusalem med Lamper, og jeg vil hjemsøge
Mændene, som sige hos sig selv: Jehova gjør
hverken godt eller ondt\dots{} Det er Lyden af Jehovas
Dag, da hviner Helten høit. En Vredesdag er denne
Dag, en Dag med Nød og Trængsel, med Bulder og
Brag, med Mulm og Mørke, med Skyer og Taage; den
Dag, da Basunen gjalder imod de befæstede Stæder og
imod de høie Taarne. Jeg vil ængste Mennesker, saa
de skulle gaae i Blinde, thi de synde mod Jehova. % PRINTED AS IS: no closing mark
(Zephan. 1).

Ingen iblandt alle Guderne har kunnet velsigne, som
Israels Gud, naar han velsigner; og Ingen iblandt alle
Guderne har kunnet forbande, som Israels Gud, naar
han forbander. Den, der vil lære, hvad religiøs
Phantasie er, han læse det Gamle Testamente. Men
Aabenbaringstanken skrider frem --- som Solen paa sin
lyse Vei til Middagshøiden, det Gamle Testamente forbereder
det Nye; vi skulle see, hvilken Forandring der
herved foregaaer med den religiøse Phantasie.

% --- p. 57 ---
\setcounter{footnote}{0}%
\lecture{Femte Forelæsning.}\opage{57}

Det kunde ved en foreløbig Betragtning næsten see ud,
som om det Nye Testamente --- naar undtages Johannes’
Aabenbaring --- ikke indrømmede den religiøse Phantasie
et saa vidt udstrakt eller storartet Spillerum som det
Gamle. I en vis Forstand maa dette vel kunne siges.
Phantasieformerne blive nemlig desto mere fremtrædende,
jo skarpere den nærværende, mangekantede Virkelighed
tegner sine Omrids mod den lyse, dæmrende Grund af
et Hiinsidigt, et Ubegrændset og Usynligt, som dog begrændser
sig og gjør sig synligt. Slagskyggerne paa
Jorden blive desto længere og mørkere, jo mere skraanende
Lyset fra oven falder paa de fremragende Høider.
Lys og Mørke, Helligt og Syndigt, Lydighed og Oprør,
et haardnakket Folk og en tugtende Jehova, disse Tilværelsens
yderste Modsætninger befinde sig paa det
gammeltestamentlige Trin under uafladelig Tiltrækning
og Frastøden i uophørlige Brydninger. Den Almægtige
selv, den uendelig Ophøiede nedlader sig til Mennesket
og taler med Mennesket: Jehova taler med Moses som
en Mand med sin Næste, og dog --- kun Jehova er
evig, Mennesket er Støv og Aske. Her er Noget, der
kan sætte den religiøse Phantasie i Bevægelse.

% --- p. 58 ---
\setcounter{footnote}{0}%
\opage{58}I det Nye Testamente  synes den religiøse Phantasie
at komme til Ro; det Overnaturlige gaaer i dybeste
Hemmelighed sammen til Eenhed med det Naturlige,
det Guddommelige bliver menneskeligt; Gud selv
aabenbarer sig som Menneske, et virkeligt Menneske,
der fødes og voxer, lider og døer. Den oprindelige
Eenhed af Gud og hans Aabenbaring er Ordet. Som
Ordet selv er den Usynlige synlig. „Og Ordet blev
Kjød, og boede iblandt os.“ Gjennem den udvortes
Ringhed lyser den indvortes Høihed, og den religiøse
Phantasie finder Hvile i Beskuelsen, naar Evangelisten
tilføier: „vi saae hans Herlighed, en Herlighed som den
Eenbaarnes af Faderen fuld af Naade og Sandhed.“
Det er en Aabenbarelse, som kan fylde den religiøse
Phantasie med Virkelighed. Dette Syn alene opveier
og forklarer alle det Gamle Testamentes prophetiske
Syner. „Salige de Øine, der see det I see!“ Men det
maa sees i Troens Lys; Vantroens Øie seer kun det
Forargelige.

Forskjellen mellem det Gamle og det Nye Testamentes
religiøse Phantasie er begrundet i Forskjellen
mellem det Gamle og det Nye Testamentes Aand.
Væsenlig er Aanden den samme, men dog paa de to
Aabenbaringstrin forskjellig. Hvor forskjellig er ikke
Loven fra Evangeliet!

Lovgivningen paa Sinai, forfærdelig for den menneskelige
Phantasie, er ledsaget af Rædsler. Bjerget bæver,
Lynene flamme; ud af Mulmet og de mørke Skyer taler
Jehova til Moses og til Moses alene; intet levende
Væsen, Menneske saa lidt som Dyr, tør nærme sig
Bjergets Fod. Virkningen beskrives i talende Træk:
„da Folket fornam Tordenen og Flammerne og Bjerget,
% --- p. 59 ---
\setcounter{footnote}{0}%
\opage{59}som røg, da flyede de og stode langt borte. Og de
sagde til Moses: tal du med os; da ville vi høre, men
lad Gud ikke tale med os, at vi ei skulle døe. --- Det er % PRINTED AS IS: no closing mark
Lovens Aand i Frygtens Phantasie; og Frygt maa der
til, Ærefrygt for det Hellige, naar et endnu raat og
uopdraget Folk skal løftes op paa Aandens Trin.

Hvor ganske anderledes med Kjærlighedsloven i
det Nye Testamente, den nye Lov, som paa een Gang
fuldkommer og ophæver den gamle. Denne Lov truer
ikke og tvinger ikke; den kjender ikke Forbandelser,
nei kun Velsignelser. Det er en Frihedens Lov, en
Hjertets Lov, thi det er en „Salighedernes“ Lov. Vi
høre det af hans Mund, der prædiker paa Bjerget.
„Salige de Fattige i Aanden, salige de Rene af Hjertet,
salige de Sagtmodige, de Barmhjertige, de Fredsommelige,
salige de, som lide Forfølgelse for Retfærdigheds Skyld,
--- Himmeriges Rige er deres.“ Denne Lovgivning indjager
ikke Skræk, dette „Salighedernes Bjerg“ er ikke
omgivet af Rædsler. Andre Syner, lyse Billeder stille
sig frem for Phantasien. Folket staaer ikke langt borte,
slaget med Forfærdelse for den Usynlige; tværtimod,
det nærmer sig i Tryghed den Synlige, det samler sig,
det flokker sig omkring ham i tætte Skarer. Ethvert
hans Ord er en Livets Kilde, Sjælene indsuge det, Øret
drikker det som en ny Viin. Intet Jordskjælv, ingen
Lyn og Tordenbrag om Bjerget, hvor han prædiker;
nei, Fred i Naturen og Fred i Menneskenes Hjerter.
Hvor er det livsaligt, foraarsfriskt, høitideligt derude:
Himlen høi og Luften klar, Fuglene frie for al Bekymring,
og Lilierne paa Marken saa tause, stille, som
lyttede de til den guddommelige Tale.

Men hvem er han nu egenlig, den nye Prophet, hvorfra
% --- p. 60 ---
\setcounter{footnote}{0}%
\opage{60}stammer han, hvis Søn er han? saa spørges der trindt
omkring i Galilæas Stæder, og Spørgsmaalet sætter
Menneskenes Phantasie i Bevægelse. Det skal være
virkeligt, at han er en Tømmermands Søn; Folk i
Nazareth kjende jo baade hans Fader og Moder, hans
Brødre og hans Søstre. --- Men han taler med en
Myndighed, som var han stegen ned fra Himlen. Hvorfra
har han denne Myndighed? Som en Prophet fra
Aanden! --- Men han taler jo som den, der staaer over
alle Propheter; hvis Søn er han? --- Skulde dette med
hans Herkomst staae hen i al Tvetydighed, saa blev
han jo en mythisk Person, men Aabenbaring er det
Modsatte af Mythe. Til Bebudelsesevangeliet slutter
sig da Fødselsevangeliet.

Det Billede, som her fremstilles for den religiøse
Phantasie, har alle Aabenbaringens Mærker. Englesynet
og det himmelske Chor er ikke bestemt for
Hyrderne alene; det er bestemt for alle Slægter, at
sige for Alle, som trods Verdens Forgudelse af overmægtige
Kraftgenier ville vende deres Tanke mod
Bethlehem, og betragte Barnet i Krybben.

I Fødselsevangeliet viser sig ret en eiendommelig
Modsætning og tillige eiendommelig Sammensmeltning
af udvortes Virkelighed og Aandsvirkelighed. Barnet i
Krybben hører jo lige saa vel som Maria, Joseph og
Hyrderne til det udvortes, historisk Virkelige; men
Englene hvor høre de hen? Til Phantasiens Verden!
Og dog er deres Tilværelse ingenlunde for den religiøse
Phantasie, hvad den er for Vantroen, kun en overtroisk
Indbildning eller i det Høieste en blot æsthetisk-mythisk
Illusion. Vexelbestemmelsen af det Sandselige og
det Oversandselige er her saa gjennemgaaende, at den,
% --- p. 61 ---
\setcounter{footnote}{0}%
\opage{61}der ikke kan see Englene, kan heller ikke see Barnet.
Barnet maa sees med det samme Troens Øie, hvormed
Englene sees; i modsat Fald et det ikke Christusbarnet, % PRINTED AS IS: „et“ for er
men et underskudt Barn, ikke det Barn, som Hyrderne
tilbede.

Den dobbelsidige Virkelighed viser os hen til den 
religiøse Phantasies guddommelige Oprindelse. Leder
man i det Udvortes efter Fødselsevangeliets historiske
Kilder, da leder man forgjæves; thi dets Oprindelse
hyller sig i et uigjennemtrængeligt Mørke. Paa den
ydre Virkeligheds Omraade ville Vantro og Overtro
herom kunne føre deres Strid indtil Verdens Ende.

For Troen er Striden afgjort; den har et Sandhedsmærke, som er ubedrageligt, og det er i Aanden. Om
Maria havde afgivet mundtlig Beretning, om Hyrderne,
saamange de vare, havde efterladt hver enkelt for sig
eller alle i Forening skriftligt Vidnesbyrd: det vilde
ikke føre os et eneste Skridt videre. Det Beviis, som
heraf kunde uddrages, var alene det, at Hyrderne og Maria havde troet paa Aabenbarelsen. Det maatte jo
da blive en Troesvirkelighed; men idet Historiens 
Virkelighed i religiøs Phantasie forvandler sig til
Troesvirkelighed, er den i samme Nu forvandlet til
Aandsvirkelighed. --- „Og strax var med Engelen en
himmelsk Hærskares Mangfoldighed, som lovede Gud“\dots{}
Den som, naar han har fæstet sit Blik paa Barnet,
ikke kan faae Syn paa den himmelske Hærskares Mangfoldighed,
han mangler religiøs Phantasie, og Folianter
af historiske Beviser ville ikke formaae at tvinge ham
til Tro. Men den, hvis Øie een Gang blev opladt til
at see --- Barnet i Krybben og Engleskaren i den
hellige Nat, for ham staaer Synet klart og vist: han
% --- p. 62 ---
\setcounter{footnote}{0}%
\opage{62}forlanger intet udvortes Vidnesbyrd. Den, som ikke
har Øre til at høre de Saliges Chor, vil lytte forgjæves
til Kritikens høirøstede Stemmer. Men den, hvis indre
Øre een Gang er blevet opladt for det himmelske Budskab,
han vil formaae at skille Sandhed fra Indbildning,
Aabenbaring fra Mythe.

Da Christus havde endt sin Bjergprædiken, for
undrede Folket sig saare over hans Lærdom, thi „han
lærte dem, som den, der havde Myndighed.“ Hvorfra
har han da denne Myndighed? Fra Himlen. Vel
sandt; men han er jo et virkeligt Menneske, født og
opvoxet som Menneske. Myndighed forudsætter personlig
Overlegenhed. Han maa altsaa under sin Forberedelse
have oplevet Noget, som paa guddommelig
og dog menneskelig Viis hører Charakteerudviklingen
til, og gjør hans uendelige Selvvished, hans indre personlige
Overlegenhed forstaaelig. Dette fører os til
Evangeliet om Fristelsen i Ørken.

Denne Fristelseshistorie danner et organisk Led i
Fremstillingen af Jesu Liv derved, at den i faa, men
store Træk giver os et Billede af den Kamp, hvori
hans besluttede Villie aabenbarer sin Reenhed og
Styrke. Men at denne Historie, saa sandt den er
virkelig, hører ind under Aandsvirkelighed, Virkelighed
for Aand, kan dog vist Ingen betvivle. Her, om ellers
nogensteds, er der Anledning til at tale om Phantasiens
Magt; thi i Fristelsen gaae Tanke og Villie
uimodsigelig sammen med Phantasie. Fristelsen er
ingen selvisk sandselig Gjæring i Christi Sind, men en
Styrkeprøve med Mørkets Aand. Her er et Sammenstød
af Villie med Villie, Christi Villie og Djævelens
Villie.

% --- p. 63 ---
\setcounter{footnote}{0}%
\opage{63}De to Villier holde sig adskilte, men den virkelige
Adskillelse, der betinger Sammenstødet, er ikke i det
udvortes Rum, den er i Phantasiens Medium. Det Udvortes
og blot Naturlige, den lange Fasten, Ørkens
Eensomhed, de øde Omgivelser, alt Sligt danner her
Relief for det Indvortes.

Det er tre Villiestanker, der alle angaae Misbrug
af de overnaturlige Gaver, hvormed en Messias, Guds
Søn, maatte være udrustet, eller rettere tre Former af
den samme Villiestanke, som Fristeren griber og anvender.
Den første er en Misbrug af Mirakelgaven,
der sigter paa det umiddelbart Sandselige: „er du
Guds Søn, da siig, at disse Stene blive til Brød“. Det
var en Phantasie, der menneskelig seet, kunde friste
den Hungrige. Havde Christus givet efter for denne
Fristelse, saa havde med det Samme hans Sendelse
været forgjæves. Den der sandselig selvisk vil bruge en
overnaturlig Gave til at stille sin Hunger, sin Tørst,
vil ogsaa sandselig selvisk bruge den til at afværge
Lidelser, Forhaanelser og trodse Spotterne paa Golgatha
ved at stige ned af Korset. Det hører med til det
Kjærlighedsoffer, som den sig til Forsagelse Indviende
bragte, at han, der ved et Mirakel kunde bespise Tusinder
i Ørken, ikke vilde røre en Steen for ved et Almagtsord
at tilfredsstille sin egen Hunger. „Da tog
Djævelen ham med sig til den hellige Stad og satte
ham paa Tindingen af Templet“. Mon denne Begivenhed
vinder i virkelig Forstaaelse ved at Troen forvandler
sig sil Overtro og gjør det Hele til en udvortes Historie? % PRINTED AS IS: „sil“ for til
Christus maatte da være ført gjennem Luften af Djævelen
omtrent som en Doctor Faust paa Mephistopheles’
Kappe. Her vindes Intet, men her tabes Alt: Vantroen
% --- p. 64 ---
\setcounter{footnote}{0}%
\opage{64}faaer Ret, naar den forvandler det Hele til en Mythe.
Det er i Aanden, baaren af en dæmonisk Phantasie,
at Christus staaer paa Templets Tinde. Den Indskydelse,
at styrte sig ned, for at lade sig bære paa Hænderne
af Engle og forvandle sin messianske Optræden til
en storartet Selvforherligelse, et guddommeligt Eventyr,
hvori der ikke behøvedes hverken Taalmod eller Lydighed
eller den Selvfornedrelsens Smerte, som skulde ende
med en Død paa Korset --- denne Indskydelse kommer
ikke fra Christi Sind, nei, den kommer fra det Væsen,
hvis Villie er Oprør mod Gud. Messias lader sig løfte
op til denne Høide af Mørkets Aand; thi han svimler
ikke paa Templets Tinde: Fristerens Forsøg er ham et
Middel til at fuldkomme sin Villies Alvor, at vie sig til
ubetinget Underkastelse under sin himmelske Faders
besluttede Raad.

At Kampen er mellem Aand og Aand, Mørkets og
Lysets, Løgnens og Sandhedens Aand, en Villieskamp i
Phantasiens Medium, bliver især tydeligt af det, som
følger. „Atter tog Djævelen ham med paa et saare
høit Bjerg og viste ham alle Verdens Riger og deres
Herlighed“. Hvor ligger vel det høie Bjerg, hvorfra
der er Udsigt over alle Verdens Riger og deres Herlighed?
I Phantasiens Land, paa Kysten ved det aabne
Hav --- i Dybet boer Forventning, og Mulighedernes
Bølger rulle --- derfra er Udsigt: paa den vilde Strøm
hist ude ligger Sirenernes Ø, stormslagne Skibe synke
som Vrag; høie Taarne sees i Luften, Slotte og Borge
i Skyernes Lag, store Stæder og vidtstrakte Egne, som
koglende Indbildning befolker. Paa Kysten hist, i
Phantasiens Land, der ligger Bjerget, det høie Bjerg
med Udsigt over al Verdens Herlighed.

% --- p. 65 ---
\setcounter{footnote}{0}%
\opage{65}Mangen En, høit begavet, rigt udrustet, mangen
Helt, hvis Navn end glimrer i Historiens Aarbøger, har
staaet paa det høie Bjerg med Mulighederne foran sig
og Verdensbetingelserne omkring sig; men han tog feil
af Skin og Virkelighed. Aandsvirkelighed var ham kun
en Phantasie; saa kom den fristende Villie hviskende,
listende frem af sit Skjul, forblommet som i en Mythe.
Den Udkaarede valgte det udvortes Store; Egoismen greb
i Virkelighedens Strenge --- et kraftigt Greb! Strengene
tonede, bruste og brast, men Grundakkorden lød paa
misbrugt Aand: den Udkaarede faldt og drog Millioner
efter sig i Faldet. Hvordan han steg, hvor høit han
stod og hvor dybt han faldt! Historien seer paa det
Virkelige, og den gigantiske Virkelighed er som en
Mythe. Saa er det da til Verdens Frelse, at En har
staaet paa det høie Bjerg, en Eneste, som med fuld
Bevidsthed mødte Mørkets Aand og gjennemskuede Bedraget.
Saa er det da til Verdens Frelse, at der trods
de høie Raab paa Villieskraft og Villiesstorhed i Selvforgudelsens
bedaarende Phantasie, endnu kan høres en
Stemme, en myndig Stemme, som lyder fra Bjerget:
„Vig bort, Satan, thi der er skrevet: du skal tilbede
Herren din Gud og tjene ham alene.“

Er det her Fremsatte nu ogsaa virkelig skikket til
at overbevise ethvert fornuftigt Menneske om Aabenbaringens,
særlig om Christendommens Sandhed?
Vi svare aabenlyst: Nei, ingenlunde! Skulde Aabenbaringens
Sandhed begrundes med almeengyldige, uimodsigelige
Fornuftbeviser, da maatte den jo begrundes
videnskabeligt. Men ingen Dannet i Nutiden vil finde
det anstødeligt, at ethvert Forsøg med at faae enten
% --- p. 66 ---
\setcounter{footnote}{0}%
\opage{66}Bebudelses- eller Fødsels eller Fristelseshistorien som
virkelig Historie sikkret enten ved historiske eller almeenphilosophiske
Grunde, fortjener at vises bort fra
de videnskabelige Begrundelsers Enemærker. Vi befinde
os med det hele Aabenbaringsindhold ingenlunde
i Videns, men i Troens Sphære. Den, der erklærer, at
Tro og Troeserkjendelse ingen Gyldighed har for ham,
at alt det om det Overnaturlige er hans Natur imod og
virker som et Tryk paa Hjertekulen, et Forspil til Noget
der begynder med Kvalme, har ikke alene Ret til at
erklære sig, men ogsaa Krav paa en vis Anerkjendelse,
for saa vidt han siger sin Mening reent
ud og uden Omsvøb melder sig som en Vantro.

Anderledes forholder det sig derimod med dem,
som i Videns og Videnskabens Navn angribe den aabenbarede
Religion. De beraabe sig ikke paa deres Vantro,
men paa deres Sandhedskjærlighed. Dersom det i
Videns klare, ubedragelige Lys kunde sees, at der var
Sandhed i Troen paa det Overnaturlige, saa vilde de
gjerne troe; men det er saa usandsynligt, og de ere saa
sandhedskjærlige! Det er enten et Selvbedrag eller et
Hykleri. Videnskaben bør skee sin Ret; men den
tør ikke overskride sin Grændse. Det er nødvendigt,
at det afgjørende Stridspunkt holdes skarpt
for Øie. Her paa Grændsen mellem Tro og Viden
spørges ikke om Sandsynlighed, her spørges alene om
Mulighed. Lad Videnskaben godtgjøre, at et Aabenbaringsphænomen er usandsynligt, høist usandsynligt,
ja, det Usandsynligste af Alt: siger Intet, slet Intet!
Men dersom dets Umulighed kan bevises, saa er Troessphæren
gjennembrudt, og dens Lys udslukkes. Den
moderne Videnskab holder Vagt ved Grændsen, og hver
% --- p. 67 ---
\setcounter{footnote}{0}%
\opage{67}Gang den gjør et Udfald, høres Feltraabet: „Umuligt!“
Lad os i Korthed see, om dette „Umuligt“ ikke skulde
være en falsk Phantasie.

Vi kaste et Blik paa Læren om Lyset. Physiken
har opstillet den høist sandsynlige Hypothese, at Lyset
beroer paa en Svingning i Ætheren. Den, der gjør sig
bekjendt med den høiere Optik, vil ogsaa i denne Henseende
have rig Anledning til at beundre den menneskelige
Skarpsindighed. Det er beviist, at det røde Lys
maa give over 500, det violette over 700 Billioner
Svingninger i et Secund. Men er nu hermed Lysets
Væsen forklaret? Ingenlunde. Lyset som Lys, som
Klarhed, forudsætter et Syn, en seende Bevidsthed, der
ikke lader sig aflede af Æthersvingninger. „Naar Ætheren
er i Hvile, er der intet Lys“; men naar alle Øine lukkes,
naar der intet seende Væsen er, hverken i Himlen
eller paa Jorden, saa er der heller intet Lys. At tale
om et Verdenslys, der er uafhængigt af alle Individers
Øine, er at tale om et Verdenssyn, der er uafhængigt
af alle seende Sjæle, d. e. at tale om Skaberens evige
Seen, om Almagtens guddommelige Phantasie.

Dersom Rummet og Ætheren og de Betingelser,
der sætte den i Svingning, ikke oprindelig fremgik af
den guddommelige Seen, saa maatte Læren om Lyset
være en Lære om et uforklarligt Skin uden Væsen.
Er der nu i den hele videnskabelige Lyslære et eneste
Støttepunkt for den Paastand, at Skaberordet: bliv Lys!
skulde høre til det Umulige? Ordet er jo netop et
kategorisk Udtryk for den guddommelige Tankes Eenhed
med den guddommelige Villie i Phantasiens altskuende
Magt. I Phantasiens Form er Ordet det
oprindelige Væsensudtryk for det Synliges Fremgang
% --- p. 68 ---
\setcounter{footnote}{0}%
\opage{68}af det Usynlige. Kan der nu i Rummet, i Ætheren,
dens Bølger og Svingninger paavises en eneste Grundbestemmelse,
som beviser Umuligheden af en altskuende
Guddoms uendelige Seen? Ingen, slet ingen. „Umuligt“
er her et tomt Udbrud af tom Indbildning, en falsk Phantasie.
Men dersom Videnskaben ikke kan bevise Skaberordets
Umulighed, naar Talen er om den Usynliges
Synliggjørelse af det naturlige Universum, hvor vil den
saa hen med sit „Umuligt“, naar Talen er om den frie
Kjærlighedsaabenbarelse i Personlighedsordet selv, det
Ord, som er „Verdens Lys“, det skabende Ord, som
var i Begyndelsen?

Indvendinger, som for en videnskabelig Kritik
ere af større Betydning, hentes fra den iøinefaldende
Modsætning mellem det, som skete paa Chisti Tid, og % PRINTED AS IS: „Chisti“
det som i vor Tid antages for muligt. „Hvem kan i
vore Dage troe paa Djævlebesættelser eller Opvækkelser
af de Døde? Kan man end, hvad Djævlebesættelser angaaer,
tye ind i den religiøse Phantasies skyggerige
Smuthuller og der søge Skjul mod den forfølgende
Kritik, saa gaaer det dog ikke an at gjøre de Dødes
Opvækkelse til en Phantasie. At opvække en Død er
lige saa umuligt som at gjøre en Steen til et Menneske.
Den, der troer paa Sligt, maa bryde Staven over alle
Videnskabens Resultater. Et af To: enten maa man
slukke Fornuftens Lys og give sig hen til den mørkeste
Overtro, eller indrømme Videnskaben med dens Inductioner,
Analogier og vidtstrakte Erfaringer Ret til at
fordre Grundsætningen om en uafbrudt sammenhængende,
jævnt fremskridende Naturudvikling respecteret, en
Grundsætning, der er af saa meget større Vægt, som
den alene stemmer med, hvad Fornuften lærer os om
% --- p. 69 ---
\setcounter{footnote}{0}%
\opage{69}det guddommelige Væsens Uforanderlighed.“ --- Saa længe
Synspunktet bliver staaende indenfor Grændsen af Sandsynligt
og Usandsynligt --- altsaa udenfor Troeserkjendelsernes
Lyskreds --- ville slige Indvendinger gjøre
Indtryk; men, naar Spørgsmaalet er om Muligt og
Umuligt, opløse de sig i et Intet. Naar Videnskaben
hævder, at den i Naturen ikke er i Stand til at opdage
Noget, der kunde tyde hen paa et Overnaturligt, og
derved tilstrækkelig godtgjør, at Læren om det Overnaturlige
ikke hører ind under det, som videnskabelig
kan begrundes: har den saa med det Samme beviist,
at det Overnaturliges Indgriben i Naturens Løb er
umulig? Naar Videnskaben hævder det guddommelige
Væsens Uforanderlighed, har den saa med det Samme
beviist, at denne Uforanderlighed nødvendig maa gjøre
den guddommelige Selvbestemmen, Guds Villies Frihed,
umulig? Naar Videnskaben hævder Umuligheden af, at
den absolute Aand kan modsige sig selv, har den saa
med det Samme beviist Umuligheden af, at denne
naturlige Tingenes Orden kan i sin usynlige, for
menneskelig Viden tilgængelige Sammenhæng med det
Overnaturlige, være Gjennemgangsled og Middel for
høiere Formaal? Har man maaksee beviist et saadant % PRINTED AS IS: „maaksee“
høiere Formaals Umulighed eller Umuligheden af, at
et Middel kan tjene sit Formaal? Naar Videnskaben
hævder --- og det kan endogsaa hævdes uden Videnskab
--- at et dødt Legeme ikke kan kaldes tillive ved
naturlige Midler, har den saa med det Samme beviist
Umuligheden af, at det kan skee ved overnaturlige
Midler? --- Vantroen mener at gjøre sig ukjendelig ved
at svøbe sig ind i Videns Kappe, men Vinden blæser
i Kappen, og paa Kløerne kjender man Løven.

% --- p. 70 ---
\setcounter{footnote}{0}%
\opage{70}Her er, vel at mærke, ingenlunde Tale om væsenlig
Splid mellem Viden og Tro. Den negative Kritik maa
have Lov til at bevise, at Aabenbaringstroen er uden
videnskabelig Gyldighed; naar den blot beviser det ret
videnskabeligt, reent videnskabeligt, saa veed man hvad
det er: det gjør slet ingen Skade. Her er alene Tale
om den Troesfornegtelse, som under Sandhedskjærlighedens
og den strenge Videnskabeligheds Skin er blevet
sig selv vigtig i Bestræbelsen med at undergrave
Religionens religiøse Grundlag.

Vantroen er driftig, snild og politisk; den forstaaer
saa godt, hvor det passer i dens Kram, at alliere sig
med Overtroen, afvise Aandsvirkeligheden og fordre det
udvortes Haandgribelige. --- „Er der,“ siger den negative
Kritik, „Noget, som Troen maa erkjende for virkeligt,
udvortes historisk Virkeligt, da er det dog nok Christi
Opstandelse.“ Ja! svarer Overtroen. „Gaaer ikke Christus
efter sin Opstandelse omkring med det selvsamme
materielle Legeme, som det, hvormed han blev begravet?“
Jo! svarer Overtroen, med det selvsamme materielle
Legeme. „Er det ikke forblommet Undvigelse og Halvtro
her at tale om Aandsvirkelighed og tye til Phantasiens
Hjælp, for at unddrage sig Tvivlens Indvendinger?“
Jo vist er det forblommet Halvtro med denne Aandsvirkelighed,
svarer Overtroen. --- Naar nu Vantroen ad
denne Vei har faaet Overtroen paa Glatiis, saa hoverer
den og slaaer det Hele i Vrangen, „Christus gjør sig
usynlig, han forsvinder for de To, som gik til Emaus;
altsaa har han intet materielt Legeme. Christus gaaer
igjennem lukkede Døre; altsaa har han intet materielt
Legeme. Christus farer op til Himlen, tvertimod Tyngdens
Lov, altsaa har han intet materielt Legeme.“ --- Den
% --- p. 71 ---
\setcounter{footnote}{0}%
\opage{71}stokhistoriske Overtro staaer betuttet, aner, at den har
drukket Dus med Skarpretteren, og veed hverken ud
eller ind. Vantroen smiler, tager den Betuttede under
Hagen, seer den ind i Øinene og siger: Virkelighed!
Ja, Du er dog et godt ærligt Stykke Virkelighed; den,
der siger om Dig, at du er en Aandsvirkelighed, Virkelighed
for Aand, han lyver Dig skammelig paa. Bliv
ikke vred; jeg kan i Grunden godt lide Dig; men ---
Du er dum.

Troen selv kan Vantroen ikke fordrage, deres
Sind og Syn er uendelig forskjelligt. Det er ikke blot
Virkeligheden, men ogsaa Phantasien, hvorom de ere
uenige. Al Troesvirkelighed er, selv om det ikke bliver
udtalt med klar Bevidsthed, altid Aandsvirkelighed;
men Aandsvirkelighed er en Forskandsning, som den
negative Kritik med al sin Snildhed, sine Reflexionsapparater
og taktiske Manøvrer ikke er istand til at
gjennembryde. Naar den mener at have faaet Tag i
det Udvortes, saa forvandler det sig og gaaer over i et
Indvortes, og naar den saa mener at kunne fastholde
det Indvortes for sig som Phantasie og Indbildning,
saa forvandler det sig igjen og slaaer ud i et Udvortes.
Utaalmodig og forbittret udbryder Kritiken: „maae
vi faae at vide, hvad og hvormeget her er Udvortes,
maae vi faae at vide, hvad og hvormeget der skal henregnes
til det Indvortes!“ Men Troen svarer uanfægtet:
her er Mysteriet, og Mysteriet er ingen Opgave for
Viden. Vantroen bilder sig ind at slaae det afgjørende Slag
og træffe Sømmet paa Hovedet, naar den sammenfatter
Undersøgelsernes Resultat i afskrækkende Slagord:
„Mystik, Mysticisme!“ Heller ikke dette gjør
Indtryk paa Troen. Mysteriet frister ikke til over%
% --- p. 72 ---
\setcounter{footnote}{0}%
\opage{72}spændt Intuition eller ufrugtbar Phantaseren; det angiver
kort, kategorisk: den menneskelige Videns Grændse.
Aabenbaringsmysteriet giver ingen Anviisning paa at
udviske Virkelighedens Mærker og liste sig uden om
Lovene; den giver tvertimod Anviisning paa en Lov, en
Frihedens Grundlov, der er lige saa sikker for Troen
som Videnskabslovene ere for Videnskaben; denne
Grundlov er en Forvandlingslov.

Forvandlingsloven kommer tilsyne i det Overnaturliges
vidunderlige Indvirken paa det Naturlige. At denne
Lov ikke tages forfængelig, er sikkret derved, at den
er en Lov for det Usandsynlige. Troen har fra sin
Side et lige saa skarpt Blik for det Usandsynlige som
Vantroen fra sin. Sløves Blikket for det Usandsynlige,
gaaer Troen øieblikkelig over til Overtro; thi Overtroen
er dum. Men Troen er ikke dum, den veed, hvad den
gjør, naar den trods Usandsynligheden griber Muligheden.
Heri ligger, at den Mulighed, Troen griber imod al
Sandsynlighed, maa for Sjælens Frelse, for Hjertets
Fred være af en saa gjennemgribende Betydning, at
den ikke kan opgives, uden at Forvisningen om det
evige Liv med det Samme opgives. Ogsaa Troen
veed, hvad den gjør: den har sine Grunde.

Vi kunne prøve Forvandlingsloven paa Opstandelsesevangeliet.
Visselig --- det er Troens Forvisning
--- er Christus opstanden fra de Døde med sit
virkelige Legeme, ikke med et Skinlegeme. Men, at dette
virkelige Legeme desuagtet hæver sig over det blot
Udvortes og Materielle, at det allerede tilhører Aandsvirkelighedens
Sphære, giver sig jo aabenbart tilkjende
ved den Opstandnes Ord til Maria Magdalene i Urtegaarden.
Det Overnaturlige virker igjennem det Naturlige.
% --- p. 73 ---
\setcounter{footnote}{0}%
\opage{73}Den Opstandnes Legeme gjennemgaaer en Række Forvandlinger.
Disse Forvandlinger, der svare til de forskjellige
Aabenbarelser for Disciplene og ende med
Himmelfarten som den sidste Forklarelse, høre væsenlig
med til Afslutningen af Jesu Liv, fordi de i hans
Personlighed gjøre det oprindelige Grundforhold mellem
det Naturlige og det Overnaturlige aabenbart.

Forvandlingsloven i og for sig er saa langt fra at
staae i Strid med Naturlærens store Grundsætninger, at
det tvertimod er den, der alene er i Stand til at begrunde
Fornuftens Ret. At forstaae Naturen uden Aand
af den bare Materie er umuligt. Naturphilosophien
leder os til det uimodsigelige Resultat, at Naturen
skabes, bæres og forklares af Aanden. Grændsen mellem
Tro og Viden er bestemt derved, at Viden maa holde
sit Resultat i Nødvendighedens Form, og kan ikke gjøre
Brug af de Frihedsmuligheder, som heri dæmre. Det
er Troens uangribelige Forudsætning, at den faste,
nærværende Tingenes Orden, hvor jævnt den end udvikler
sig, hvor mange Millioner- Billioner Aar dens
Udviklingsgang end udviser, dog har sin inderste
oprindelige Grund i et Anlæg af skabende Frihed; at,
med andre Ord, det hele Universum hviler paa Grundlaget
af et skjult Under.

Aabenbaringen gjør det Skjulte aabenbart. Opstandelsesmiraklet
anskueliggjør Forholdet mellem det
Forkrænkelige og det Uforkrænkelige, mellem Her og
Hisset, mellem Døden og Livet. Om Aabenbaringsmysteriet
veed Physiken med al sin Mathematik Intet,
slet Intet. Den har i „sit Auditorium“ Ret til at tale
om Vandets og Legemernes Vægtfylde, om Tyngdelovens
Almeengyldighed i den naturlige Verden o. desl.; men
% --- p. 74 ---
\setcounter{footnote}{0}%
\opage{74}naar Vantroen i den Anledning gjør sig et Ærende ind
i Aabenbaringens Lyskreds, for i Tyngdelovens og Vægtfyldens
Navn at raabe: Umuligt! forraader den en Aandsraahed,
der maa finde sig i at blive viist udenfor.

Det skal villig indrømmes, at den aabenbarede
Religions Indhold har endnu mangfoldige andre Sider,
hvorfra den maa belyses end netop Phantasiens. Desuagtet
turde det vel være paa Tide at blive opmærksom
paa, at den store Forvirring paa Troens Omraade for
en stor Deel hidrører fra en Forkvakling af den religiøse
Phantasie. Medens der i alle Dannelses Kredse
omhyggelig sørges for at den æsthetiske Phantasie med
Sands og Smag for Alt, hvad der henhører under Konst
og Poesie, bliver behørig udviklet, er den religiøse Phantasie
traadt i Skygge, endog i den Grad, at dens Uddannelse
synes at være afsluttet med Confirmationsalderen.
En Blanding af alskens Elementer, lidt fra
Æsthetik og lidt fra Physik, lidt fra Mythologie og lidt
fra videnskabelig Kritik forvirrer Synet og fordunkler
Opfattelsen af Aabenbaringsordet saavel i det Nye som
i det Gamle Testamente. Men Phantasien er en Magt,
ogsaa i Religionen; Enhver, som føler Trang til at
gjøre sig Rede for sin Tro, vil sikkerlig ogsaa føle
Trang til at klare sig sin religiøse Phantasie.

% --- p. 75 ---
\setcounter{footnote}{0}%
\lecture{Sjette Forelæsning.}\opage{75}

Hvad der i en Henseende er begrændset, kan i en
anden Henseende være ubegrændset; dette gjælder
øiensynligt om Phantasiens uendelige Rige. At dette
Rige er uendeligt, maae selv de indrømme, der fornegte
dets guddommelige Væsen og søge Phantasiens Udspring
i en physiologisk Disposition af Nerver og Hjerne.

Der er i enhver til Bevidsthedslivet hørende Virksomhed
en skjult Phantasie. Ingen Hukommelse er
mulig uden Phantasie; Erfaringsbillederne, Billeder af
hvad vi have oplevet, af hvad vi have seet i Naturen
eller lært i Historien, ligge ikke hengjemte i Sjælen
som i et Magasin: de ere i vort Indre som Muligheder,
som Dispositioner i vort Nervesystem. At vi mindes
det Forbigangne, beroer paa, at Phantasien ifølge en
Villiestilskyndelse kan paa vilkaarlig Maade vække de
indsovede Muligheder og stille det Oplevede frem i et
Billede for vort indre Syn. Ogsaa Forstanden med
dens sindige lovbundne Virksomhed er betinget af
Phantasie. Det er Forstanden, der adskiller, sammenligner,
ordner og forbinder vore Forestillinger; men
hvad er en Forestilling? Et sig opløsende Billede, svævende
mellem det Usynlige og det Synlige, det Ubestemte
og dog Bestemte: som Billedets Opløsning er, maa
% --- p. 76 ---
\setcounter{footnote}{0}%
\opage{76}ogsaa Billedet være, men Billedet dannes af Phantasien.
Mathematiken er i strengeste Forstand en Forstandsvidenskab;
men Mathematik er umulig uden tilsvarende
Phantasie. De utallige Combinationer af Former og
Formler, de mangfoldige sammensatte, fordoblede, ofte
høist indviklede geometriske Figurer gjøre det indlysende
nok, at der maa gives en egen mathematisk Phantasie.
Heller ikke Villien, hvor praktisk den er og hvor udelukkende
den end tager Sigte paa det Virkelige, kan
undvære Phantasiens Bistand. Den Opgave, som Villien
sætter sig, er jo endnu ikke virkelig løst, men den skal
løses, det vil sige; den er en Mulighed, hvis Indhold
maa være anskueligt i Phantasiens Form. Og hvor
inderlig er ikke Phantasien forenet med Følelsen; man
tænke paa Stemning, Glædens Stemning, Sorgens Stemning:
ingen ideel Stemning uden Phantasie!

Men Phantasien, hvor meget den end er til Tjeneste
for enhver Sjælevirksomhed ved enhver Anledning, nøies
dog ikke med en simpel Tjenerrolle; den har en
hemmelig, man kunde sige kvindelig Drift til selv at
herske. Den kommer ikke plumpt med det; den gaaer
varlig frem for Trylleriets Skyld: den bruger List. Den
melder sig ikke som en stortalende Herre, der træder
i Gulvet med klirrende Sporer og fører Kommando-Ordet;
nei, som en væver lille Pige, med Skudsmaal:
ferm til sin Gjerning og villig til Alt.“ Men just for
denne Fermheds og Villigheds Skyld maa man holde
Øie med den, bestandig have den i Kikkerten og nøie
vogte paa hvert dens Skridt. Den seer saa uskyldig
ud og er dog saa udspekuleret, listig, intrigant, en
Skuespillerinde af første Rang, med Mesterskab i Forvandlinger,
Trylleoptog og bedaarende Illusioner. Phan%
% --- p. 77 ---
\setcounter{footnote}{0}%
\opage{77}tasiens Rolleregister er af et uhyre Omfang fra det
Dybsindigste og Rørende indtil det Lystigste, meest
Overgivne, Taabelige og Fjantede, der er til.

Ihvorvel Phantasien er sindig, bøielig, bevægelig og
kan optage Alt og bearbeide, hvad den optager paa
den allerfrieste Maade, saa har den dog et forskjelligt
Præg hos de forskjellige Individualiteter. Man seer det
paa det Eiendommelige i enhver Art af Begavelse. Der
gives en egen Malerphantasie, en musikalsk Phantasie,
Digterphantasie o. s. v. Hvor Phantasiebegavelsen i en
eller anden Retning overskrider sin Grændse, begynder
øieblikkelig det Søgte og Maniererede. Ja ikke blot
hos de særlig Begavede, men hos ethvert velorganiseret
Individ har Phantasien sin egen Grundform. Det er
forstaaeligt nok, naar vi betænke, at den Anskuelsens
Eenhed, hvori alle Livsindtryk, Tanker, Forestillinger,
Forbilleder og Bestræbelser optages, netop afgives af
Phantasien. Ligesom ethvert Menneske i legemlig Forstand
har sit eiendommelige Ansigt, saa har det ogsaa
i sjælelig Forstand sit eiendommelige, af dets Grundphantasie
bestemte Physiognomie: Phantasien er Sjælens
ideelle Ansigt. Uden et høiere Vexelskin af de to
Ansigter, det indre og det ydre, vilde det Mimiske i
Skuespillet tabe sin Betydning. Til at være en stor
Mimiker hører noget Mere end at kunne lægge sit
Ansigt i Folder og skjære Grimacer. Al plastisk Inspiration
gaaer igjennem Phantasien.

Dette er saa meget naturligere, som Phantasien
ogsaa omfatter den Side af Bevidstheden, der igjennem
det Halvbevidste rækker ned i det Ubevidstes dunkle
Regioner. Phantasien arbeider lige saa vel i Bevidsthedens
Nat som i dens Daglys; den elsker Dæmringen,
% --- p. 78 ---
\setcounter{footnote}{0}%
\opage{78}Skumringen og den vaagne Drøm, hvori dens luftige
Skuespil allerhelst opføres. Hermed stemmer det, at
Phantasiens Magt ikke umiddelbart afhænger af bevidste
Tanker eller et vilkaarligt Forestillingsløb, men gjennemtrænger
Kjød og Blod og virker som en aandig Naturkraft.
Man tænke paa den Angst og Skræk som et
Blik eller et Ord kan vække, man tænke paa den uvilkaarlige
Blegnen og den lige saa uvilkaarlige Rødmen.

Derfor bør man aldrig maale et Menneskes personlige
Dannelse med hans Forstandsudvikling, hans Indsigter
og Kundskaber alene; ville vi have en Maalestok
for hans Sindsstyrke, maae vi have Phantasien med,
den Elementærphantasie, som boer i Nerverne. Er
denne Phantasie forkvaklet eller forsømt, saa har Villien
intet sikkert Grundlag og Besindelsen kan svigte En i
det afgjørende Øieblik. Samling, Fatning, Aandsnærværelse
beroer væsenlig paa umiddelbar Phantasie.
Et raat Menneske, hvis Phantasie er lige saa afstumpet
som hans Samvittighed, kan i visse Tilfælde udvise
større Koldblodighed end en Dannet, hvis Elementærphantasie
ikke er blevet udviklet i Forhold til hans
Forstandsoplysning. En Kirketyv, der bryder ind ved
Nattetid, kan famle omkring ved Alteret og mellem
Kisterne i en Begravelse med en Rolighed saa uanfægtet,
som om han gik i sin Dagligstue. Griber Nogen ham
pludselig i Mørket bagfra, saa falder det ham dog ikke
ind, at det er et Spøgelse; han veed strax, hvad det er:
det er Politiet. Derimod kan mangen Oplyst holde sig
fuldstændig overbeviist om, at der ingen Gjengangere
gives, og dog ved Nattetid være bange for at søge Selskab
med de Døde: han er egenlig ikke bange for
de Døde, men han er bange for sin Phantasie.
% --- p. 79 ---
\setcounter{footnote}{0}%
\opage{79}Den bekjendte La Mettrie, en af de allerkrasseste
Materialister mod Slutningen af forrige Aarhundrede, en
Mand, der hverken var bange for Gud eller Djævelen,
var bange for at gaae over en Kirkegaard om Natten.

Den, der har en syg Phantasie, kan umulig have
en sund Sjæl. Som man taler om Forhærdelse f. Ex. i
Leveren, saaledes kan man ogsaa tale om Forhærdelsesknuder
i Phantasien; hvad ere vel de saakaldte Fixideer
Andet end saadanne Knuder. Alle Sjælesygdomme,
hvilke de organiske Aarsager end kunne være, vise hen
til en Forstyrrelse i Phantasien. Vel er Forstanden
ogsaa forstyrret; men at Forstandsforvirringen hidrører
fra forstyrrende Billeder, vilde Indskydelser, kan blandt
Andet sees deraf, at de Gale kunne ræsonnere med
skarp Forstand og sikker Logik indenfor de Grændser,
som Fixideerne afsætte. Hvor uhyggeligt at gjøre et
Indblik i et Galehuus: alle Phantasiens Forvrængninger
blive der synlige.

Der er et Billede af Kaulbach, det forestiller en
Daarekiste. Man seer de Sindssyge samlede og grupperede,
men man seer ogsaa, at den Ene slet ikke
ændser den Anden: Enhver af dem er sig selv nok,
udelukkende beskjæftiget med sin egen Fixidee. De
forskjelligste Indbildninger og Lidenskaber staae malede
i Ansigter og Stillinger. Især er der En, --- Figuren
er ingen Karikatur --- men han opvækker Gysen. Det
er en Begeistret, der forkynder Evangeliet; han løfter
Korset og smiler som en Forklaret. Men i dette Blik,
i dette salige Smiil, hvilket bundløst Dyb af Vanvid!

Phantasien er en væsenlig Form for Sandhedens
Indhold og dog første Opfinder af alskens Indbildninger,
Løgn og Bedrag. Fra psychologisk Side ligger da
% --- p. 80 ---
\setcounter{footnote}{0}%
\opage{80}heri en Opfordring til nøiere at undersøge Phantasiens
Oprindelse og gjøre sig Rede for, hvorfra den faaer sit
høist uligeartede Indhold. Psychologien maa i vore
Dage støtte sig til Physiologien; uden Indsigt i Organbygningen
og de organiske Betingelser er det umuligt
at komme til klar Erkjendelse eller frugtbar Indsigt i
de sjælelige Virksomheder. Ogsaa i denne Henseende
har den fremskridende Naturlære vundet en betydelig
Indflydelse paa Aandslæren. Herhen hører en nærmere
Opklaring af de saakaldte Hallucinationer. Det er
bekjendt hvorledes den Febersyge kan phantasere,
hvorledes Drankeren i sit Delirium faaer Syner af en
saa haandgribelig Realitet, at han tumler med dem,
som med legemlige Gjenstande. Opiumsspiserens indbildte
Herligheder og deres gradvise Forvandling til
lige saa mange virkelige Elendigheder, høre til de
Phænomener, der vække Opsigt.

I J. L. Heibergs samlede Skrifter (8de Bind) findes
en lille Afhandling om en engelsk Opiumsspisers Bekjendelser.“
„Opium,“ hedder det, „bringer Sjælen i
en rolig, ja alvorlig Nydelsestilstand. Alle Sjæleevnerne
synes at udvikle sig til en overordenlig Grad, idet de
dog tillige vedligeholde den fuldkomneste Harmonie.“
Vi faae her Exempler paa, hvad en materiel Paavirkning
af Nervesystemet har at betyde. „Jeg stod,“ siger
Bekjenderen, „høit over Livets Tummel; Forvirring,
Bekymring, Ængstelse, de hemmelige Byrder, som
undertrykke Hjertet, alt var forsvundet; det forekom
mig som en Helligaften for hele Skabningen, som en
Hviledag for alle Menneskehedens Elendigheder. Det
jordiske Livs Glæder og Forhaabninger vare forsonede
med Graven, og min aandelige Bevægelse foregik uden
% --- p. 81 ---
\setcounter{footnote}{0}%
\opage{81}Anstrængelse, uden Træthed, ligesom Stjernernes paa
Himlen. Og med alt dette en dyb Ro, som ikke fremkom
af en Uvirksomhedstilstand, men syntes at være
Resultatet af uendelige Kræfters Modsætning, det var
paa een Gang en Virksomhed uden Grændse og en
uendelig Hvile“. Med hvilken Magie en sygelig anspændt
Phantasie kan virke, i hvilken Grad den formaaer at
fremkogle en Verden af Syner og give dem Præg af
Virkelighed, fremlyser af følgende Beskrivelse: „Henflyttet
under Vendekredsens Himmel, saae jeg mig“ ---
man erindre sig, at Manden er i London --- „snart omgiven
af alle disse Egnes overordenlige, baade vegetabile
og animale Frembringelser. Fugle, Abekatte,
Reptilier omringede mig fra alle Sider, og forfulgte
mig uophørlig. Jeg flygtede ind i Pagoder og fandt
mig indelukket for evig i disse Bygningers hemmelige
Smuthuller. Jeg var paa een Gang Afgudsbilledet og
Præsten, Guden og Offeret. Jeg flygtede gjennem Asiens
Skove, forfulgt af Bramas Vrede; Vischnu hadede mig,
og Siva lagde Snarer for mig. Paa een Gang stod jeg
lige for Isis og Osiris, de beskyldte mig for at have
begaaet en Forbrydelse, hvorover Ibis og Krokodillen
skjælvede og bævede. De dømte mig til en tusindaarig
Nat i Gravmæler af Steen, og begravede mig med
Sphinxer og Mumier i utilgængelige Pyramider. Derpaa
fandt jeg mig nedsænket i Nilens Mudder, midt imellem
Uhyrer, gjennemtrukne af Floddynd, og mellem alle
Slags rædsomme og væmmelige Ting“. --- Billederne
blive bestandig mørkere og vildere. „En uhyre Krig,
et Slag, en Dødskamp maatte indtræffe paa et eller
andet Sted, og oprulle sig for mig som et mørkt og
sørgeligt Skuespil; men jeg vidste hverken hvor eller
% --- p. 82 ---
\setcounter{footnote}{0}%
\opage{82}hvorledes, eller ved hvem. Den levende Deeltagelse,
jeg følte derfor, var saa meget utaaleligere, som jeg
hverken kjendte Beskaffenheden eller Aarsagen eller
det sandsynlige Udfald af det, som jeg interesserede
mig for. Som det ofte hænder os i Drømme, hvor vi
stedse ere Midtpunktet af alt det, som tildrager sig
omkring os, saaledes forekom det mig ogsaa her, at
jeg havde Magt til at bestemme dette Udfald, og dog
tillige, at jeg ikke havde det. Jeg følte, at jeg kunde
gjøre det, dersom jeg havde Kraft til at ville det, men
jeg var ligesom nedtrykt af en uhyre Byrde, som jeg
ikke kunde befrie mig fra. Jeg forblev altsaa i min
Uvirksomhed. Men nu nærmede sig det kritiske Øieblik,
Oprøret blev stærkere, et uendeligt Anliggende
skulde afgjøres. Derpaa fulgte en pludselig Bestyrtelse,
overilet Flugt i alle Retninger, og Gjenlyd af utallige
Flygtningers Skridt. Derpaa gik Skygger og Blus,
Hvirvelvinde og menneskelige Ansigter mig forbi; og
endelig --- under Følelsen af, at Alt var tabt --- Fruentimmerskikkelser,
og Trækkene af det, som havde været
mig kjærest i Verden. Og kun et eneste Øieblik var
tilstaaet os; vi maatte skilles for evig. Alle vred Hænderne,
jamrede og hylede, og fra alle Sider hørte man
gjentage: Farvel for evig! --- Disse rædselsfulde Ord gjenlød
i mit Indre, jeg vaagnede i den skrækkeligste Ængstelse,
og skreg med høi Røst: Jeg vil aldrig mere sove!“

At Sligt maa hidrøre fra Affectioner i Hjernen,
derom kan der ikke være Tvivl. Physiologien har i
denne Henseende stor Fortjeneste af at have gjort
Rede for de forskjellige Arter af Hallucinationer.

Dog, ogsaa i dette Punkt maa Aandslæren være
paa sin Post, thi Sagen kan let faae en Vending, som
% --- p. 83 ---
\setcounter{footnote}{0}%
\opage{83}forvansker Opfattelsen, nedsætter det Ideale til Skin
og forvrænger det Allerhøieste.

Ligesom Kritiken ved Hjælp af en dybere gaaende
Psychologie har opdaget Mythernes Udspring og forvandlet
Aabenbaring til Mythe, saaledes kan den nu,
overfladisk seet, med en vis Ret beraabe sig paa en
dybere gaaende Physiologie og afvise de historiske
Vidnesbyrd, der bære Præg af psychologisk Troværdighed,
hvor Talen er, ikke om dunkle, halvbevidste,
men om klare, fuldt bevidste Aabenbarelser, ved den
korte og tørre Bemærkning: det er Hallucinationer.

Et Par Exempler. Den hellige Antonius flygter
fra Verden ud i Ørken, men Djævelen følger med og
plager ham med allehaande Fristelser. Snart er det
hæslige Trolde, der skrække ham, og snart deilige
Kvinder, der sætte hans Afholdenhed paa Prøve; man
hører langt borte hans Skrig i Hulen, thi Satan selv
angriber ham corporlig og slaaer ham til Jorden. Heri,
siger Physiologen, er intet Urimeligt. Det kan Altsammen
virkelig være gaaet saaledes til. Eremiten har
plaget sig med Askese, Fasten og Spægelser saalænge,
indtil han er blevet et Offer for de forfærdeligste
Hallucinationer.

Det fortælles om Muhamed, at da han i Eensomheden
under Bøn og Fasten grublede paa sin nye Religion, da hørte han fra alle Sider i den omgivende
Natur Stemmer, som stadfæstede hans guddommelige
Sendelse. Det talede fra Klipperne: Du er Profeten;
det lød fra Skyerne: Du er Propheten; Træerne stak
Hovederne sammen, Løvet hviskede og Kilderne rislede:
See der gaaer Profeten. Og Muhamed blev ført
gjennem
alle syv Himle. I den øverste saae han Dødens
Engel;
% --- p. 84 ---
\setcounter{footnote}{0}%
\opage{84}ja, han saae Guds største Skabning, en Engel, som
havde 70,000 Hoveder, i hvert Hoved 70,000 Munde, i
hver Mund 70,000 Tunger; med hver Tunge talte den
70,000 Tungemaal og forkyndte i dem alle Herrens
Priis. Tilsidst blev han ført frem for Gud Herren selv,
og til Høire for Thronen stod der at læse: der er ingen
anden Gud end Gud, og Muhamed er hans Profet. Da
Koreishiternes Stamme satte sig op imod ham og ikke
vilde lytte til hans Prædiken, fordi han krænkede
deres Forret ved at kæmpe mod Afguderiet, var hans
Farbroder Abu-Taleb den Eneste, der tog ham i
Forsvar, skjøndt han negtede at bekjende hans Tro.
Han raadede sin Brodersøn til at afstaae fra sit Forehavende,
da deraf ellers vilde flyde store Ulykker; men
Muhamed svarede: „Kom man ogsaa til mig med Solen
i den ene Haand og Maanen i den anden, man skulde
dog ei faae mig til at vige“.

Hvad er nu Indbildning og hvad er Sandhed? Kan
dette i Troens Sphære afgjøres med objektivt videnskabelige
Grunde og Beviser? Umuligt! Havde vi end
i hver Stad 70,000 theologiske Faculteter, og hvert Facultet
70,000 lærde Theologer, og hver Theolog 70,000
Folianter, og hver Foliant 70,000 theologiske Beviser:
det skulde dog ikke rokke en Muhameds Tro, skulde,
selv om han havde forstaaet hvert eneste Beviis til
Punkt og Prik, ei faae ham til at vige et Haarsbred.
Kan dette saa maaskee gjøre os usikkre over Troeserkjendelsens
Gyldighed? Aldeles ikke. Synets Sands
kan bedrage, hvi skulde Phantasien ikke ogsaa kunne
bedrage. Men saa lidt som noget fornuftigt Menneske
vil opgive at stole paa sit Syn, uagtet han veed, at
Synet kan skuffe, lige saa lidt vil nogen troende Sjæl
% --- p. 85 ---
\setcounter{footnote}{0}%
\opage{85}forarges over, at Aabenbaringen kommer i Phantasiens
Form, ihvorvel han veed, at Phantasien kan skuffe.

For at skjelne skarpt imellem Sandt og Falsk,
maae vi bestandig have Opmærksomheden fæstet paa
Phantasiens væsenlige Indhold i dets Eenhed med
Formen. Dette gjælder ikke alene om det Religiøse,
hvis Indhold er Aandsvirkelighed, Virkelighed for Aand;
men det gjælder ogsaa i Sphærer, hvor Indholdet er, hvad
vi særlig kalde Idee. Det gjælder i Konst og Poesie,
det gjælder praktisk i den rationelle Ethik. Konsten
har sin Sandhed, uagtet der gives Forvrængninger af
Konst; Poesien har sin Sandhed, og denne Sandhed er
netop kjendelig paa, at der gives en usand Poesie, en
tom og fad, ja usædelig og for de Ubefæstede høist
fordærvelig Poesie.

Tør Poesien være usædelig? Der er ført megen
Strid om Forholdet mellem det Skjønne og det Gode.
De, som fordre, at Poesien skal være moralsk i den
Forstand, at den umiddelbart skal moralisere, de glemme,
at der hvad Indholdsformen angaaer, gives to væsenlig
forskjellige Ideer, den æsthetiske Idee, det Skjønne, og
den ethiske Idee, det Gode; de glemme, at Phantasien
er Selvhersker i Skjønhedens Rige og har Ret til at
tage Alt i sin Tjeneste, ikke blot det Gode, men ogsaa
det Onde, det Umoralske, det Hæslige. Men de, der
gaae til den modsatte Yderlighed og paastaae, at Poesien
slet ikke behøver at ændse de sædelige Fordringer,
de, der antage, at den frie Phantasie ligerviis som den
frie Tanke kan i ethisk Forstand byde Læseren hvad
det skal være; at den med Virkeligheden og dens Muligheder
for Øie, netop feirer sin meest glimrende
Triumph, naar den forstaaer at gribe en Bordelhistorie,
% --- p. 86 ---
\setcounter{footnote}{0}%
\opage{86}et Emne i sine prosaiske Omrids saa gemeent, at modne
Mænd for Velanstændighedens Skyld med Uvillie vise
det fra sig, og dog udstyre dette lurvede Emne saa
glimrende, saa brillant og pikant og lade Fremstillingens
Trylleri forherlige det Uanstændige endog i den
Grad, at unge Læserinder, uden at vide, hvad de gjøre,
fordybe sig deri og draabeviis faae deres Phantasie
forgiftet; --- de, der opfatte det Poetiske fra dette
Synspunkt, maae paa Forhaand have lukket deres Øine
for Poesiens guddommelige Mission.

Det er ikke blot, hvad man fra et høitflyvende
Emancipationens Standpunkt maaskee vil kalde en
spidsborgerlig Moral, som herved bliver haanet: det er
Ideen selv, der er krænket. Man gjøre Kløften mellem
det Æsthetiske og det Ethiske saa bred, man vil, saa
er det dog vel aabenbart, at den ethiske Idee har en
idetmindste lige saa gyldig Realitet som den æsthetiske.

At der lige saa vel maa gives en ethisk, som en
æsthetisk Phantasie, er let at oplyse. Men i det Ethiske
er Phantasien tjenende. Det er allerede bemærket, at
Villien trænger til Phantasiens Hjælp. Intet er farligere
for den sædelige Villie end en forvirret, forhexet, af
falske Billeder vildledt Phantasie. Hvorfra stammer
Forfængelighed? Fra falsk Phantasie! Og Ærgjærrighed
og Pengegridskhed? Fra forhexet Phantasie! Og Vellystens
Dæmon? Fra forgiftet Phantasie! Den, der
planmæssig lægger an paa at fordærve et Menneske,
vil ikke indskrænke sig til at omstøde disse eller hine
Grundsætninger, eller lokke den Paagjældende ind i
disse eller hine Vaner; nei! Fordærveren kan ikke være
sikker paa sit Offer, før han i Bund og Grund har
faaet dets Phantasie fordærvet. Først med den for%
% --- p. 87 ---
\setcounter{footnote}{0}%
\opage{87}dærvede Phantasie ville de sidste Mindelser af Samvittighedsnag
høre op, og den sidste Rest af en sædelig
Idee være med Rode oprykket.

Men det er dog vel indlysende, at de sande Ideer,
hvor forskjellige de end indbyrdes ere, umulig kunne
være hverandres Fjender og nære en saadan Ringeagt,
ja et saa indgroet Had til hverandre, at den ene skulde
lægge an paa at forhaane eller endog udrydde den
anden. Forholdet er jo lige det omvendte: der er i
Grunden god Forstaaelse mellem Idee og Idee, der er
Fred og Harmonie i Ideernes Sphære.

Vel kan der ogsaa tales om en Ideernes Egoisme,
idet en enkelt Idee kan trænge saa stærkt paa, at den
skyder de andre tilside; men denne Egoisme hører ikke
til Ideernes indre Væsen; den kommer ved deres Gjennemgang
gjennem det menneskelige Sind, paa Overgangen
fra Evighed til Tid, til Udtræden i Verden, til Virkeliggjørelse
i det Timelige. Det er med Ideen i Menneskets
Sind som med Blommen til Fuglens Æg, først er den
lille og blød som et Bær, dernæst som en Drue og tilsidst
ved Gjennemgangen gjennem Æggelederen overtrækkes
den med Hinden og den haarde Skal. Og dog
kommer Ideen anderledes frem af sit Anlæg end Fuglen
af Ægget. Jo mere en Idee er beaandet af det Almene
og gjennemtrængt af Aand, desto klarere vil den under
sin individualiserende Fremtræden vise os, at den med
al sin Eiendommelighed vedkjender sig
Slægtskabet
med de øvrige Ideer, især med dem af den nærmeste
Kreds, hvorunder den hører. Dette gjælder bogstavelig
om Forholdet mellem den æsthetiske og den ethiske Idee.

Den „moderne Poesie“ gjør vel i at
tage „Verdens
Børn“ som de ere, og blotte den Huulhed og Indbildskhed,
% --- p. 88 ---
\setcounter{footnote}{0}%
\opage{88}hvormed en afbleget Fromhedsphantasie eller reglementeret
Kasernemoral saa gjerne vil føre det store
Ord; den gjør vel i at indsætte det naturlige Menneskes
Drifter og Anlæg i deres oprindelige Rettigheder.
Men disse Drifter og Anlæg maae jo dog have et
menneskeligt Maal og være Aandens Love undergivne.
Det er et uhyggeligt Tegn paa forvirret Phantasie, naar
end ikke Digteren kan faae Øie paa Aandslovene, men
gjør Driftens Tilfredsstillelse selv til Ideal, gjør sine
Helte og Heltinder til Meteorer og Sæbebobler, og smiler
ironisk, naar de stærke Sjæle, forviklede i deres egne
Conseqvenser, gaae tilgrunde. Man kan ikke bebreide
den moderne Poesie, at den fordyber sig i det Virkelige
og vinder i Mesterskab, hvor det gjælder om at assimilere
sig dets Elementer i levende Former og charakteersvangre
Situationer. Men man kan med Grund bebreide
den, at dens Livsbilleder ere falske, fordi de have
enten en fad og tom eller endog aldeles falsk Baggrund.
Det Poetiske lider ikke derved, at det nærværende,
virkelige Liv med dets Modsætninger af selvvigtige
Borneertheder og ekstastiske Emancipationer faaer Lov % PRINTED AS IS: „ekstastiske“
at træde frem i tydelige Træk; men det lider forfærdeligt
derved, at Indtrykket af den hele Fremstilling viser,
at Digteren selv med al sin Begavelse er fattig paa
Aand og kun driver Leg med Ideer, et verdensklogt
og dog i sandselig Emancipation forvirret Subject, der
mangler Maalestok. Den moderne Poesie --- Cherbuliez’s
Miss Rowel f. Ex. --- forstaaer meget godt at forlige
det Ethiske med det Æsthetiske, for saa vidt den har
Noget deraf, men den har ikke Meget.

Efter det evige Væsen maa det Gode absolut være
skjønt, og det Skjønne godt. Ikke desto mindre er
% --- p. 89 ---
\setcounter{footnote}{0}%
\opage{89}Phantasieformen for disse Ideers Fremtræden høist
forskjellig.

I Konst og Poesie bearbeider Phantasien det Virkelige
alene fra Mulighedssiden; det vil sige: den faaer
kun Virkeligheden med i Speilbilleder af det Virkelige,
men det Virkelige selv rører den ikke med en Finger.
Den lader den reelle Verden gaae sin egen skjæve Gang
og bygger sig til Erstatning en ideal Verden, rigere,
herligere og mere fortryllende end den, som er dannet
af det grove Stof. Og Menneskene forlystes ved Synet
af den nyskabte Verden; Lys og Skygge hvor ere de
lykkelig fordeelte, Lysbrydningerne mangfoldige og Farveblandingen
smeltende, yndig, henrivende. Her blotter
Naturen sine dulgte Skjønheder, sine jomfruelige Hemmeligheder.
Her gaaer det med Fart over Bjerge og Dale;
selv Modgangen er et Krydderi, Sinkerierne, belles
horreurs, ere interessante. Og Menneskene i denne
Verden, hvor ere de aabenhjertige! De betroe os Alt,
hvad de føle og tænke. Man kommer dem saa nær,
man kan fornemme deres Aandedrag og mærke, hvor
Pulsen slaaer. Her er Befordring, her er Erfaring, her
vindes Menneskekundskab: man har „Mosters Briller“
paa, og seer hver Enkelt lige ind i Hjertet.

Med den ethiske Phantasie er det anderledes tarveligt. Idet Phantasien tager Tjeneste hos Villien, maa
Virkeligheden jo vistnok ogsaa opfattes fra Mulighedernes
Side; men med de ethiske Muligheder forholder det sig
paa en ganske anden Maade end med de æsthetiske.
De eneste Muligheder, som Villien ændser, er dem, som
virkelig lade sig realisere. Villien seer med Phantasiens
Blik paa Muligheden, men, vel at mærke, for Virkelighedens
Skyld; den ethiske Villie gaaer fra Virkelighed
% --- p. 90 ---
\setcounter{footnote}{0}%
\opage{90}gjennem Mulighed over til en ny Virkelighed, thi den
omarbeider Virkeligheden.

Er nu den nye, af Villien bearbeidede Virkelighed
skjønnere og rigere end den, hvormed der begyndtes?
Vi svare: Ja! kun at Skjønheden ikke væsenlig er i det
Udvortes, men i det Indvortes. Lægges det afgjørende
Eftertryk paa Forskjønnelsen af det Udvortes, saa tager
det Æsthetiske Magten fra det Ethiske, og den hele
Bestræbelse, om der end arbeides i Ansigtets Sved,
faaer et overveiende Præg af blot æsthetisk Phantasie.

Her er et Phantasiebedrag, for hvilket Culturverdenen
med sin rastløse Arbeiden for Livets Bekvemmeligheder,
Comfort i det Mindre og storartede
Forskjønnelser i det Større, store Anlæg, store Bygninger,
store Mindesmærker o. s. v., idelig er udsat.
Bedraget bestaaer ingenlunde deri, at der arbeides i
det Materielle paa Fremgang og Forskjønnelse, langtfra!
Bedraget er, at denne Arbeiden i udadgaaende Retning
sættes som et usvigeligt Tegn paa det Ethiske. Bedraget
er, at det Ethiske, langt fra at holde Skridt
med Kulturen, tvertimod svækkes, idet Pragtphænomenernes
Glands giver Maalestokken an og sætter det
Godes usynlige Skjønhed i Skygge.

Det Godes Skjønhed er, naar vi lægge Vægt paa
Modsætningen til det Æsthetiske, en usynlig Skjønhed.
Vel har denne usynlige Skjønhed ogsaa sit Phænomen,
og kan sees i Phantasie; men den æsthetiske Betragtning
giver kun en forbigaaende Stemning: den sande
Indsigt er betinget af personlig Oplevelse. Trofasthed,
Standhaftighed, Taalmodighed, Langmodighed, ere usynlige
Skjønheder. Det er ikke de ydre Omgivelser, som
her skulle glimre; tvertimod! selv den æsthetiske Phan%
% --- p. 91 ---
\setcounter{footnote}{0}%
\opage{91}tasie forstaaer, at de --- for Effectens Skyld --- helst
maae være saa tarvelige som muligt. Et mørkt Fængsel,
en fattig Stue, et Sygeleie, det er Omgivelser, som
passe for den usynlige Skjønhed.

Den æsthetiske Phantasie kan digte Taalmodigheden
og forelske sig i sit Digt; men den virkelige Taalmodigheds
Væsen fatter den ikke. Den kan forelske
sig i et skjønt Ansigt; men naar den forelsker sig i
Taalmodighedens Skjønhed, falder der et Slør over dens
Øine, thi den seer ikke Væsenet selv, men kun et Skin,
og begriber ikke, hvorledes den taalmodige Sjæl i
Virkeligheden kan holde ud at være saa taalmodig.
Dette Spørgsmaal kan alene den ethiske, ja strengt
taget den religiøs-ethiske, Phantasie forstaae, men ---
paa Oplevelsens Betingelse.

Vi begyndte med at see os om i Phantasiens store
Rige, og nu vi have seet, at dette Rige er uendeligt,
er vor korte Vandring til Ende. Vil Nogen sige ved
sig selv: det er Skade, thi nu var det just værd at
begynde, da vilde det være mig ikke blot en Paamindelse,
at denne Fremstilling visselig har sin Ufuldkommenhed,
men ogsaa et Vidnesbyrd om, at vi dog
paa vor korte Vei have mødt Adskilligt, hvormed det
vel var værd at gjøre nærmere Bekjendtskab.

\end{document}
